\documentclass[Total.tex]{subfiles}

\begin{document}
	\tableofcontents
	\chapter{Projective and Injective Modules, Free Modules}
		\section{Free Modules}
			\begin{Def}
				
			\end{Def}
			\subsubsection{Universal Properties}
			\begin{proposition}
				Suppose $R$ be a ring, $M$ be $R$-Mod and $I\in {\mathcal Set}$ be index, then we have a correspondence
				\begin{gather*}
					Hom_{Mod-R}(R^{\oplus I},M)\simeq Hom_{Set}(I,M)
				\end{gather*} 
			\end{proposition}
			\begin{proof}
				内容...
			\end{proof}
		\section{Projective Modules}
			\begin{Def}
				An object $P$ of a category $C$ is a \textbf{projective object} if it has the \textbf{left lifting property} against epimorphisms. 
				
				This means, that $P$ is projective if for any morphism $f:P\rightarrow B$ and any epimorphism $q:A\rightarrow B$, $f$ factors through
				\begin{center}
					\begin{tikzcd}
						&A\arrow[d,"\forall q\in Epi"]\\
						P\arrow[r,"\forall f"]\arrow[ru,"\exists",dashrightarrow]&B
					\end{tikzcd}
				\end{center}
			\end{Def}
			An equivalent way to say to say this is that
			\begin{Def}
				An object $P$ is projective precisely if the $hom-functor$ $Hom(P,-)$ preserves epimorphisms.
			\end{Def}
			
			With property of epimorphism, $P\rightarrow A$ is unique up to isomorphism.
			
			\begin{proposition}
				The followings are equivalent in $\mathcal Set$:
				\begin{itemize}
					\item Every object is projective;
					\item The axiom of choice holds.
				\end{itemize}
			\end{proposition}
			\begin{proof}
				$\Rightarrow$ In ZF, $AC\Leftrightarrow$ Every surjection $f:A\xrightarrow B$ has a section structure
				\begin{gather*}
					g\circ f=id_A
				\end{gather*} 
				We take $f=id_B,P=B$.
				
				$\Leftarrow$ If $AC$ holds, then we could give a function $P\rightarrow A$. Now we associate $f^{-1}(b)$ with $\bar{b}\in f^{-1}(b)$ by AC. Thus we could select $g:P\rightarrow \{\bar{b}|b\in B\}$ such that $f(x)=q(g(b))$.
				\begin{center}
					\begin{tikzcd}
						&C\arrow[d,"\subseteq"]\\
						&A\arrow[d,"q"]\\
						P\arrow[ruu,dashrightarrow]\arrow[ru,dashrightarrow]\arrow[r,"f"]&B
					\end{tikzcd}
				\end{center}
			\end{proof}
			
			\textbf{Remark} If we replace categorey of sets by another base topolos will behave essentially precisely like set, but in general will not validate the axiom of choice. 
			
			Homological algebra in such a more general context is abelian sheaves, and the theory is the theory of cohomology of abelian sheaf.
			
			A main task is the resolution of $R$-mod.
			
			\textbf{Remark} Although mostly, we have a forgetful functor $\mathcal{F}\rightarrow Sets$, but the functor is not full, so not preverse the functions. 
			\begin{Def}
				The \textbf{projetive module} over $R$ is the projective objects in $R-Mod$.
			\end{Def}
			\begin{proposition}
				\label{Homological-Algebra-Projective-Injective-Modules-Equivalent-Definition-of-Projective-Modules}
				The following properties are equivalent:
				
				\begin{itemize}
					\item $P$ is projective module;
					\item If $g:B\rightarrow C$ is surjective, then
					\begin{gather*}
						g_*:Hom(P,B)\rightarrow Hom(P,C) 
					\end{gather*}
					is also surjective;
					\item For short exact sequence \begin{tikzcd}
						0\arrow[r]&A\arrow[r]&B\arrow[r]&C\arrow[r]&0
					\end{tikzcd}
					Then, 
					\begin{center}
						\begin{tikzcd}
							0\arrow[r]&Hom_R(P,A)\arrow[r]&Hom(P,B)\arrow[r]&Hom(P,C)\arrow[r]&0
						\end{tikzcd}
					\end{center}
					is exact;
					\item Short exact sequence $0\rightarrow P\rightarrow F\rightarrow V\rightarrow 0$ is split exact;
					\item $P$ is a direct item of free modules, i.e. there exist some $V$ such that
					\begin{gather*}
						F=P\oplus V
					\end{gather*}
					where $F$ is free.
				\end{itemize}
			\end{proposition}
			\begin{proof}
				\begin{itemize}
					\item (2)$\Leftrightarrow$ (3)$\Leftrightarrow$ (4) By the discussion of exact sequence
					\item (3)$\Leftrightarrow$ (4) By the discussion of exact sequence.
					\item (1)$\Leftrightarrow$ (4) Consider the diagram
					\begin{center}
						\begin{tikzcd}
							&&A\arrow[d,twoheadrightarrow]\\
							F=P\oplus V\arrow[rr,dashrightarrow,bend right=20]\arrow[rru,dashrightarrow,"Epi\Rightarrow \exists p "]&P\arrow[r]\arrow[l,"\iota"]\arrow[ru,"p|P"]&B
						\end{tikzcd}
					\end{center}
				\end{itemize}
			\end{proof}
			
			
			\begin{proposition}
				Suppose $P$ be a finite projective $R$-module. Then, there exists finite free module $F$, and finite module $Q$ such that
				\begin{gather*}
					F=P\oplus V
				\end{gather*}
			\end{proposition}
			\begin{proof}
				$\Leftarrow$
				
				$\Rightarrow$ Consider the generators of $P$.
			\end{proof}
			\subsubsection{Flatness of Projective Module}
			
			\begin{proposition}
				Every projective module is flat.
			\end{proposition}
			\begin{proof}
				Conisder the alternative definition
			\end{proof}
			\subsubsection{Free Modules are Projective}
			\begin{proposition}
				Let $R$ be a ring.Every free $R$-module is projective.
			\end{proposition}
			\begin{proof}
				Consider diagram
				\begin{center}
					\begin{tikzcd}
						内容...
					\end{tikzcd}
				\end{center}
			\end{proof}
			\subsection{Sum}
			\begin{lemma}
				内容...
			\end{lemma}
			\begin{proof}
				内容...
			\end{proof}
			
			\begin{proposition}
				(Criterion of Projective Property for Sum) Suppose $P=\sum P_i$, then
				\begin{center}
					$P$ is projective $\Leftrightarrow$ $P_i$ is projective.
				\end{center}
			\end{proposition}
			\begin{proof}
				内容...
			\end{proof}
			
			\begin{corollary}
				The direct sum of projective modules is projective.
			\end{corollary}
			\begin{proof}
				Omitted.
			\end{proof}
			\subsection{Tensor}
			\begin{proposition}
				\begin{itemize}
					\item Suppose $P$ be a projective right $R$-module, then: $P\otimes_R-:{}_R M\rightarrow Ab$ is an exact functor.
					\item Similarly for $M_R$.
				\end{itemize}
			\end{proposition}
			\begin{proof}
				内容...
			\end{proof}
			
			\subsubsection{Flatness}
				\begin{Def} Let $R$ be a ring:
					\begin{itemize}
						\item (Flat) $R$-module is \textbf{flat} if $N_1\rightarrow N_2\rightarrow N_3$ is exact$\Rightarrow M\otimes_R N_1\rightarrow M\otimes_R N_2\rightarrow M\otimes_R N_3$ is exact;
						\item (Faithfully Flat) Instead of $\Rightarrow$ by $\Leftrightarrow$;
						\item $R\rightarrow S$ is \textbf{flat/faithfully flat} if $S$ is flat/faithfully flat module over $R$;
					\end{itemize}
				\end{Def}
				
				
				\begin{lemma}
					Let $R$ be a ring, $I,J\subseteq R$ be ideals. Let $M$ be a flat $R$-module. Then
					\begin{gather*}
						IM\cap JM=(I\cap J) M
					\end{gather*}
				\end{lemma}
				\begin{proof}
					Consider the exact sequence
					\begin{gather*}
						0\rightarrow I\cap J\rightarrow R\rightarrow R/I\oplus R/J\qquad\mbox{Iso III}
					\end{gather*}
					
					Tensoring with the flat module $M$, we obtain an exact sequence
					\begin{gather*}
						0\rightarrow(I\cap J)\otimes_R M\rightarrow M\rightarrow M/IM\oplus M/JM
					\end{gather*}
					
					We have
					\begin{gather*}
						(I\cap J)M=Ker(M\rightarrow M/IM\oplus M/JM)=IM\cap JM
					\end{gather*}
					
					By the definition of morphism.
				\end{proof}
				
				\begin{lemma}
					Let $R$ be a ring, let $\{M_i,\varphi_{ii'}\}$ be a directed system of flat $R$-modules. Then
					\begin{gather*}
						colim_i M_i
					\end{gather*}
					
					is a flat $R$-module.
				\end{lemma}
				\begin{proof}
					Firstly: $\otimes$ commutes with colimits;
					
					secondly: directed colimits are exact.
					
					Thus
					\begin{gather*}
						0\rightarrow A\rightarrow B\rightarrow C\rightarrow 0\Rightarrow 0\rightarrow colim_i M_i\otimes A\rightarrow colim_i M_i\otimes B\rightarrow colim_i M_i\otimes C\rightarrow 0
					\end{gather*}
				\end{proof}
				
				\begin{lemma}
					The composition of (faithfully) flat ring maps is (faithfully) flat.
					
					If $R\rightarrow R'$ is (faithfully) flat, and $M'$ is a (faithfully) flat $R'$-module, then $M'$ is a (faithfully) flat $R$-module.
				\end{lemma}
				\begin{proof}
					\begin{itemize}
						\item Consider
						\begin{gather*}
							N_1\rightarrow N_2\rightarrow N_3
						\end{gather*}
						
						an exact complex of $R$-module. Then
						\begin{gather*}
							R'\otimes_R N_1\rightarrow R'\otimes_R N_2\rightarrow R'\otimes_R N_3
						\end{gather*}
						
						Furthermore
						\begin{gather*}
							M'\otimes_{R'}(R'\otimes_R N_1)\rightarrow M'\otimes_{R'}(R'\otimes_R N_2)\rightarrow M'\otimes_{R'}(R'\otimes_r N_3)
						\end{gather*}
						
						And we have
						\begin{gather*}
							M'\otimes_{R'}(R'\otimes_R N)\cong (M'\otimes_{R'} R')\otimes_R N=M'\otimes_R N
						\end{gather*}
						\item Omitted.
					\end{itemize}
				\end{proof}
				
				\begin{lemma}
					Let $M$ be an $R$-module. The following are equivalent:
					\begin{itemize}
						\item $M$ is flat over $R$;
						\item For every injection of $R$-modules $N\subseteq N'$ the map
						\begin{gather*}
							N\otimes_R M\rightarrow N'\otimes_R M
						\end{gather*}
						
						is injective.
						\item for every ideal $I\subseteq R$, the map
						\begin{gather*}
							I\otimes_R M\rightarrow R\otimes_R M=M
						\end{gather*}
						
						is injective;
						\item For every finitely generated ideal $I\subseteq R$, the map
						\begin{gather*}
							I\otimes_R M\rightarrow R\otimes_R M=M
						\end{gather*}
						
						is injective.
					\end{itemize}
				\end{lemma}
				\begin{proof}
					\begin{itemize}
						\item (1)$\Rightarrow$ (2) By definition;
						\item (2)$\Rightarrow$ (3) We take $I,R$ as $R$-module.
						\item (3)$\Rightarrow$ (4) Trivial;
						\item (4)$\Rightarrow$ (1) 
					\end{itemize}
				\end{proof}
				
				\subsubsection{Directed Limit}
				\begin{lemma}
					Let $\{R_i,\varphi_{ii'}\}$ be a system of rings over the directed set $I$. Let
					\begin{gather*}
						R=colim_i R_i
					\end{gather*}
					\begin{itemize}
						\item If $M$ is an $R$-module such that $M$ is \textbf{flat} as an $R_i$-module for all $i$. Then, $M$ is flat as an $R$-module;
						\item For $i\in I$, let $M_i$ be a flat $R_i$-moudle, and $i'\geq i$ and let $f_{ii'}:M_i\rightarrow M_{i'}$ be a $\varphi_{ii'}$-linear map such that
						\begin{gather*}
							f_{i'i''}\circ f_{ii'}=f_{ii''}
						\end{gather*}
						
						Then, $M=colim_{i\in I} M_i$ is a flat $R$-module.
					\end{itemize}
				\end{lemma}
				\begin{proof}
					内容...
				\end{proof}
				
				\begin{lemma}
					Suppose that $M$ is (faithfully) flat over $R$, and $R\rightarrow R'$ is a ring map. Then
					\begin{gather*}
						M\otimes_R R'
					\end{gather*}
					
					is (faithfully) flat over $R'$.
				\end{lemma}
				\begin{proof}
					内容...
				\end{proof}
				
				\begin{lemma}
					Suppose that $R\rightarrow R'$ be a faithfuly flat ring map, $M$ be a module over $R$. Now, we set
					\begin{gather*}
						M'=S'\otimes_S M
					\end{gather*}
					\begin{itemize}
						\item 
						\item 
					\end{itemize}
				\end{lemma}
				\begin{proof}
					内容...
				\end{proof}
				\subsubsection{Via (faithfully) Flat over $R'$}
				\begin{lemma}
					Suppose that $M$ is \underline{(faithfully) flat} over $R$, $R\rightarrow R'$ is a ring map. Then
					\begin{gather*}
						M\otimes_R R'
					\end{gather*}
					
					is \textbf{(faithfully) flat over $R'$}.
				\end{lemma}
				\begin{proof}
					内容...
				\end{proof}
				
				\begin{lemma}
					Let $R\rightarrow R'$ be a \underline{faithfully flat} ring map. Let $M$ be flat over $R$ iff $M'$ is flat over $R$.
				\end{lemma}
				\begin{proof}
					
				\end{proof}
				
				\begin{lemma}
					Let $R$ be a ring, $S\rightarrow S'$ be a flat map of $R$-algebras. Let $M$ be a module oveer $S$. Set
					\begin{gather*}
						M'=S'\otimes_S M
					\end{gather*}
					\begin{itemize}
						\item 
						\item 
					\end{itemize}
				\end{lemma}
				\begin{proof}
					内容...
				\end{proof}
				
				\begin{lemma}
					内容...
				\end{lemma}
				\begin{proof}
					内容...
				\end{proof}
				\subsubsection{Equational Criterion of Flatness}
				\begin{lemma}
					A module $M$ over $R$ is \textbf{flat} iff every relation in $M$ is trivial.
				\end{lemma}
				\begin{proof}
					内容...
				\end{proof}
				\subsubsection{Flatness of Items in Exact Sequence}
				\begin{lemma}
					内容...
				\end{lemma}
				\begin{proof}
					内容...
				\end{proof}
				\subsubsection{Localization}
				\subsubsection{Directed System of Faitfully Flat Algebras}	
				\begin{proposition}
					Let $R$ be a ring, let $\{S_i,\varphi_{ii'}\}$ be a directed system of faifully flat $R$-algebras. Then
					\begin{gather*}
						S=colim_i S_i
					\end{gather*}
					
					is a faithfully flat $R$-algebra.
				\end{proposition}
				\begin{proof}
					
				\end{proof}
			\subsection{Projective Basis Theorem}
				\begin{theorem}
					Let $R$ be a ring, $P$ is an $R$-module. The followings are equivalent:
					\begin{itemize}
						\item $P$ is finitely generated projective module;
						\item Existence of projective basis: There exists a family of elements $\{x_i\}_{i=1}^n\subseteq P$, with a set of homomorphsims $\{\phi_i\}_{i=1}^n Hom_R(P,R)$, where $\phi_i$ is linear! and for all $x\in P$,
						\begin{gather*}
							x=\sum \phi_i(x)x_i
						\end{gather*}
					\end{itemize}
				\end{theorem}
				\begin{proof}
					$\Rightarrow$ Suppose $P$ is projecive, then $P\oplus Q=(R)^n$. Then, there exists:
					\begin{gather*}
						i:P\rightarrow R^n\qquad p:R^n\rightarrow P
					\end{gather*}
					
					such that $p\circ i=id_P$. Then, $x=p(i(x))=\sum a_j(x)p(e_j)$.
					
					$\Leftarrow$ Suppose there exists such basis, then define
					\begin{gather*}
						\psi:P\rightarrow R^n\qquad x\mapsto (\phi_1(x),\dots,\phi_n(x))\\
						\theta:R^n\rightarrow P\qquad (a_1,\dots,a_n)\mapsto \sum a_ix_i\\
						\Rightarrow (\theta\circ \psi)(x)=x
					\end{gather*}
				\end{proof}
				
				\textbf{Note: Flatness}
				
				\textbf{Note: Serre-Swan}
			\subsection{Projective Resolution}
				\begin{Def}
					Let $M$ be $R$-module, then a \textbf{left resolution} is a complex $P_i=0,i<0$ with a map $\epsilon:P_0\rightarrow M$ and the complex
					\begin{gather*}
						\dots\rightarrow P_1\rightarrow P_0\rightarrow M\rightarrow 0
					\end{gather*}
					is exact. An resolution is \textbf{projective} iff $P_i$ is projective.
				\end{Def}
				
				\begin{proposition}
					Every $R$-module $M$ has a projective resolution.
					
					More generally, if an abelian category $\mathcal{A}$ has enough projectives, then every object $M$ in $\mathcal{A}$ has a projective resolution.
				\end{proposition}
				\begin{proof}
					Consider the following constructions: By selecting $M_{i-1}:=ker(P_i\rightarrow P_{i-1})$ with new projective resolution $P_{i+1}\rightarrow M_i$. Now consider the compose 
					\begin{gather*}
						P_{i+1}\rightarrow M_i\rightarrow P_i
					\end{gather*}
					\begin{center}
						\begin{tikzcd}
							P_1\arrow[rr]\arrow[rd]&&P_0\arrow[r]&M\arrow[r]&0\\
							&M_0\arrow[rd]\arrow[ru]\\
							0\arrow[ru]&&0
						\end{tikzcd}
					\end{center}
					
					And $P_{n+2}$ is projective. Thus 
					
					
					
				\end{proof}
				
				\subsubsection{Comparision Theorem/ Homological Algebra's Fundamental Theorem}
				\begin{theorem}
					Let $P\stackrel{\epsilon}{\rightarrow} M$ be a projective resolution of $M$ and $f':M\rightarrow N$ be a map in $\mathcal{A}$. Then, for every resolution $Q\stackrel{\eta}{\rightarrow} N$ there is a chain map $f:P\rightarrow Q$ lifting $f'$ in the sense of
					\begin{gather*}
						\eta\circ f_0=f'\circ\epsilon
					\end{gather*}
					\begin{center}
						\begin{tikzcd}
							\dots\arrow[r]&P_1\arrow[r]\arrow[d,dashrightarrow]&P_0\arrow[r,"\epsilon"]\arrow[d,dashrightarrow,"f_0"]&M\arrow[r]\arrow[d,"f'"]&0\\
							\dots\arrow[r]&Q_1\arrow[r]&Q_0\arrow[r,"\eta"]&N\arrow[r]&0
						\end{tikzcd}
					\end{center}
					
					Moreover, such $f$ is unique up to \textbf{chain homotopy}.
				\end{theorem}
				\begin{proof}
					The prove the existence of such chain $P$ by induction:
					
					Basis Step: Consider the following square
					\begin{center}
						\begin{tikzcd}
							P_0\arrow[r,"\epsilon"]\arrow[d,dashrightarrow,"f_0"]&M\arrow[d,"f'"]\\
							Q_0\arrow[r,"\eta"]&N
						\end{tikzcd}
					\end{center}
					
					$\epsilon$ and $\eta$ are epi. So the triangle \begin{center}
						\begin{tikzcd}
							&Q_0\arrow[d,"\eta"]\\
							P_0\arrow[r,"f\circ \epsilon "]\arrow[ru,dashrightarrow]&N
						\end{tikzcd}
					\end{center}
					
					Induction Step: Consider
					\begin{center}
						\begin{tikzcd}
							P_{n+2}\arrow[r]\arrow[d,dashrightarrow,"f_{n+2}"]&P_{n+1}\arrow[r]\arrow[d,"f_{n+1}"]&P_n\arrow[d,"f_n"]\\
							Q_{n+2}\arrow[r]&Q_{n+1}\arrow[r]&Q_n
						\end{tikzcd}
					\end{center}
					
					We use diagram chasing: Consider $p\in P_{n+2}$, given a $f_{n+2}$ such that
					\begin{gather*}
						d_{n+2}\circ f_{n+2}=f_{n+1}\circ d_{n+1}
					\end{gather*}
					
					Notice that $Im(f_{n+1}\circ d_{n+2})\subseteq Ker(d_{n+1})=Im(d_{n+2})$. Thus we have the following triangle
					\begin{center}
						\begin{tikzcd}
								&Q_{n+2}\arrow[d]\\
							P_{n+2}\arrow[ru,dashrightarrow,"f_{n+2}"]\arrow[r]&Im(d_{n+2})
						\end{tikzcd}
					\end{center}
				\end{proof}
				\subsubsection{Horseshow Lemma}
				\begin{theorem}
					Suppose given a commutative diagram
					\begin{center}
						\begin{tikzcd}[cramped,sep=small]
							&&&0\arrow[d]\\
							\dots P_2'\arrow[r,"P_2"]&P_1'\arrow[r]&P_0'\arrow[r]&A'\arrow[r]\arrow[d]&0\\
							&&&A\arrow[d]\\
							\dots P_2'\arrow[r]&P_1'\arrow[r]&P_0'\arrow[r]&A'\arrow[r]\arrow[d]&0\\
							&&&0
						\end{tikzcd}
					\end{center}
					where the column is exact and the row are projective resolutions. Then, set 
					\begin{gather*}
						P_n:=P_n'\oplus P_n''
					\end{gather*}
					The right hands column lift an exact sequence of complexes
					\begin{gather*}
						0\rightarrow P'\stackrel{i}{\rightarrow}P\stackrel{\pi}{\rightarrow} P''\rightarrow0
					\end{gather*}
				\end{theorem}
				\begin{proof}
					内容...
				\end{proof}
			\subsection{Projective Module over PID}
				\begin{proposition}
					If a ring $R$ is a PID,(especially, $R=\ZZ$) then, every projecive $R$-module is free.
				\end{proposition}
				\begin{proof}
					Notice that: Every submodule of module over PID is free. Then, by the equivalent definition of projective module.
				\end{proof}
			\subsection{Kaplansky's Theorem}
				\begin{lemma}
					Suppose $P$ be a projective module, $A\subseteq P$ be a submodule. If $A$ is finitely/countably generated, then there exists an item $Q$ such that:
					\begin{itemize}
						\item 
						\item 
						\item 
					\end{itemize}
				\end{lemma}
				\begin{proof}
					内容...
				\end{proof}
				\begin{theorem}
					(Kaplansky) Every projective module is the direct sum of the finitely generated projective modules.
				\end{theorem}
				\begin{proof}
					Firstly, consider the chain of submodules:
					\begin{gather*}
						P_0\subseteq P_1\subseteq \dots \subseteq P_i\subseteq \dots
					\end{gather*}
					
					where $P_{i+1}/P_i$ is a 
				\end{proof}
				
				\begin{theorem}
					Let $R$ be a ring. Then the following are equivalent:
					\begin{itemize}
						\item $R$ is a local ring;
						\item Every projective module over $R$ is free, and has \textbf{indecomposable decomposition}
						\begin{gather*}
							M=\oplus M_i
						\end{gather*}
						
						such that for each maximal direct summand $L$ of $M$, there exists a decomposition
						\begin{gather*}
							内容...
						\end{gather*}
					\end{itemize}
				\end{theorem}
		\section{Locally Free/Finite Locally Free Modules}
			\begin{Def}
				Let $R$ be a ring and $M$ be an $R$-module:
				\begin{itemize}
					\item We say that $M$ is \textbf{locally free} if we can cover $Spec(R)$ by \textbf{standard opens} $D(f_i),i\in I$ such that: $M_{f_i}$ is a free $R_{f_i}$-module for all $i\in I$. i.e.
					\begin{gather*}
						M_{f_i}=R_{f_i}^{\oplus n_i}
					\end{gather*}
					\item We say that $M$ is \textbf{finite locally free} if we \underline{can} choose the covering such that $M_{f_i}$ is finite free.
					\item $M$ is \textbf{finite locally free} of \textbf{rank} $r$ if
					\begin{gather*}
						M_{f_i}\cong R_{f_i}^{r_i}
					\end{gather*}
				\end{itemize}
			\end{Def}
			
			\begin{lemma}
				Let $R$ be a ring and let $M$ be an $R$-module. The following are equivalent:
				\begin{enumerate}
					\item $M$ is finitely presented, and $R$-flat;
					\item $M$ is finite projective;
					\item $M$ is a direct summand of a finite free $R$-module;
					\item $M$ is finitely presented and for all $p\in Spec(R)$, the localization $M_p$ is free;
					\item $M$ is finite and locally free;
					\item $M$ is finitely presented and for all maximal ideals, $m\subseteq R$ the localization $M_m$ is free;
					\item $M$ is finite, for every prime $p$, the module $M_p$ is free, and the function
					\begin{gather*}
						\rho_M:Spec(R)\rightarrow \ZZ\qquad p\mapsto dim_{\kappa(p)} M\otimes_R \kappa(p)
					\end{gather*}
					is \textbf{locally constant} in Zaraski topology.
				\end{enumerate}
			\end{lemma}
			\begin{proof}
				We use this way to prove:
				\begin{center}
					\begin{tikzcd}
						(2)\arrow[r,Leftrightarrow]&(3)\arrow[r,Rightarrow]&(1)\arrow[r,Rightarrow]&(7)\arrow[r,Leftrightarrow]\arrow[d,Rightarrow]\arrow[rd,Rightarrow]&(6)\\
						&&(5)\arrow[u,Rightarrow]&(4)\arrow[l,Rightarrow]&(8)
					\end{tikzcd}
				\end{center}
				\begin{itemize}
					\item $(2)\Leftrightarrow (3)$: Consider the equivalent definition of projective module: $P$ is projective $\Leftrightarrow F=P\oplus V$, notice that, by the construction of $F$, such $F$ is also finite.
					\item $(3)\Rightarrow (1)$:If $R^{\oplus n}=M\oplus N$, and $n$ is finite. We have
					\begin{gather*}
						M=R^{\oplus n}/N
					\end{gather*}
					which is finitely presented. And $M$ is projective $\Rightarrow -\otimes M$ is flat.  
					\item $(1)\Rightarrow (7)$: Pick a prime $p;x_1,\dots,x_r\in M$ which map t a basis of $M\otimes_R \kappa(p)$.
					\item (6)$\Leftrightarrow$ (7)
					\item (7)$\Rightarrow$ (8)
					\item (7)$\Rightarrow$ (4)
					\item (4)$\Rightarrow$ (5)
					\item 
				\end{itemize}
			\end{proof}
			
			\textbf{Hint} $R$-flat+finite$\not\Rightarrow$ Projective.
			
			Counterexample: With $R=C^\infty(\RR),M=R_m=R/I$ where $m=\{f\in R|f(0)=0\}$ and $I=\{f\in R|\exists\epsilon,\epsilon>0:f(x)=0,\forall x,|x|<\epsilon\}$. 
			
			\begin{proposition}
				If $R$ be a reduced ring, and let $M$ be an $R$-module. Then, equivalent to:
				\begin{itemize}
					\item 
				\end{itemize}
			\end{proposition}
			\begin{proof}
				内容...
			\end{proof}
			\begin{proposition}
				Suppose: $R$ is a local ring, $M$ is a finite flat $R$-module. Then $M$ is finite free.
			\end{proposition}
			\begin{proof}
				内容...
			\end{proof}
			\begin{proposition}
				Let $R\rightarrow S$ be a flat local homomorphism of local rings. Let $M$ be a finite $R$-module. Then, $M$ is finite projective over $R$ iff
			\end{proposition}
			\begin{proof}
				内容...
			\end{proof}
				\begin{proposition}
					For any $R$-module $P$ the following statements are equivalent:
					\begin{itemize}
						\item $P$ is finite locally free
						\item 
						\item 
						\item 
					\end{itemize}
				\end{proposition}
			\subsection{Faithfully flat descent for projectivity of modules}
				\textbf{Question} Suppose $\phi:R\rightarrow S$ be a ring homomorphism, $M$ is $R$-module. Could we state, that $M\otimes_R S$ is projective $S$-module$\Rightarrow$ $M$ is projective $R$-module?
				
				In general not, but
				\begin{proposition}
					If $\phi$ is \textbf{faithfully flat},
					\begin{gather*}
						(0\rightarrow M'\rightarrow M\rightarrow M''\rightarrow 0)\in Exact(R-Mod)\Leftrightarrow  (0\rightarrow M'\otimes S\rightarrow M
						\otimes S\rightarrow M''\otimes S\rightarrow 0)\in Exact(R-Mod)
					\end{gather*} then the statement holds.
				\end{proposition}
				\begin{proof}
					内容...
				\end{proof}
		\section{Injective Modules}
			\begin{Def}
				The \textbf{injective object} in $\mathcal{A}$ satisfies the universal lifting property: Given an injection $f:A\rightarrow B$ and a map $\alpha:A\rightarrow I$ there exists at least one map $\beta:B\rightarrow I$ such that $\alpha=\beta\circ f$
				\begin{center}
					\begin{tikzcd}
						0\arrow[r]&A\arrow[r,"f"]\arrow[d,"\alpha"]&B\arrow[ld,"\beta"]\\
						&I
					\end{tikzcd}
				\end{center}
			\end{Def}
			
			We say $\mathcal{A}$ has enough injectives if for every object $A$ in $\mathcal{A}$, there is an injection $A\rightarrow I$ with $I$ injective. 
			
			\begin{proposition}
				In category $Sets$, every nonempty object is injective.
			\end{proposition}
			\begin{proof}
				内容...
			
			\end{proof}
			\subsubsection{Product}
			\begin{lemma}
				If $I_\alpha$ is a family of injectives, so is $\prod I_\alpha$.
			\end{lemma}
			\begin{proof}
				内容...
			\end{proof}
			
			\subsection{Baer's Criterion}
			\begin{theorem}
				(Baer) A right $R$-module $E\in Mod-R$ is injective iff for every right ideal $J$ of $R$, every map $J\rightarrow E$ can be extended to a map
				\begin{gather*}
					R\rightarrow E
				\end{gather*}
			\end{theorem}
			
			\begin{corollary}
				内容...
			\end{corollary}
			\begin{proof}
				内容...
			\end{proof}
			
			\subsubsection{Abelian Groups}
			
			\begin{corollary}
				$\mathcal Ab$ has enough injectives.
			\end{corollary}
			\begin{proof}
				内容...
			\end{proof}
			\subsection{Embedding Theorem}
			\subsection{Injective Resolution}
			\begin{Def}
				内容...
			\end{Def}
			
			\subsubsection{Comparison Theorem}
			
			\subsubsection{Adjoint Functors}
			
			\begin{lemma}
				For every right $R$-module $M$< the natural map
				\begin{gather*}
					\tau:Hom_{Ab}(M,A)\rightarrow Hom_{mod-R}(M,Hom_{Ab}(R,A))
				\end{gather*}
				is an isomorphism, where
			\end{lemma}
			\begin{proof}
				内容...
			\end{proof}
			\begin{Def}
				A pair of functions $L:\mathcal{A}\rightarrow \mathcal{B}$ and $R:\mathcal{B}\rightarrow \mathcal{A}$ are adjoint if there is a natural bijection for all $A\in \mathcal{A}$ and $B\in \mathcal{B}$ such that
				\begin{gather*}
					\tau:Hom_\mathcal{B}(L(A),B)\stackrel{\cong}{\rightarrow} Hom_\mathcal{A}(A,R(B))
				\end{gather*}
				Natural: For $f:A\rightarrow A'$ and $g:B\rightarrow B'$,
				\begin{center}
					\begin{tikzcd}
						Hom_\mathcal{B}(L(A'),B)&Hom_\mathcal{B}(L(A),B)&Hom_\mathcal{B}(L(A),B')\\
						Hom_\mathcal{A}(\mathcal{A}',R(B))&Hom_\mathcal{A}(A,R(B))&Hom_\mathcal{A}(A,R(B'))
					\end{tikzcd}
				\end{center}	
			\end{Def}
			\begin{corollary}
				内容...
			\end{corollary}
			\begin{proof}
				内容...
			\end{proof}
			
			\begin{proposition}
				
			\end{proposition}
			\begin{proof}
				内容...
			\end{proof}
			
			\begin{corollary}
				内容...
			\end{corollary}
			\begin{proof}
				内容...
			\end{proof}
		\section{Projective and Injective Objects in Abelian Category}
			\subsection{Projective Object}
				\begin{Def}
					Object $X\in \mathcal{C}$ is \textbf{projective} if it satisfies the left lifting property for \textbf{epimorphism}.
					\begin{center}
						\begin{tikzcd}
							&I\arrow[d,twoheadrightarrow]\\
							A\arrow[r]\arrow[ru,dashrightarrow]&B
						\end{tikzcd}
					\end{center}
				\end{Def}
				\begin{proposition}
					Let $X\in \mathcal{A},\mathcal{A}$ be an abelian category. The followings are equivalent:
					\begin{itemize}
						\item $X$ is projective.
						\item $Hom(X,-)$ is exact.
					\end{itemize}
				\end{proposition}
				\begin{proof}
					\begin{itemize}
						\item $\Rightarrow$ Consider the exact sequence
						\begin{center}
							\begin{tikzcd}
								0\arrow[r]&C_1\arrow[r,"f"]&C_2\arrow[r,"g"]&C_3\arrow[r]&0
							\end{tikzcd}
						\end{center}
						Then the morphisms induced by the morphisms above
						\begin{center}
							\begin{tikzcd}
								0\arrow[r]&Hom(X,C_1)\arrow[r,"f^*"]&Hom(X,C_2)\arrow[r,"g^*"]&Hom(X,C_3)\arrow[r]&0
							\end{tikzcd}
						\end{center}
						Injective easy to verify, and surjective given by the diagram
						\begin{center}
							\begin{tikzcd}
								&C_2\arrow[d,"g"]\\
								X\arrow[ru,dashrightarrow]\arrow[r]&C_3
							\end{tikzcd}
						\end{center}
						
						$Im(f)=Ker(g)$ from exactness of $Hom(X,-)$.
						\item $\Leftrightarrow$ If $Hom(X,-)$ is exact, now consider the following exact sequnece
						\begin{gather*}
							0\rightarrow Ker(g)\rightarrow A\rightarrow B\rightarrow 0
						\end{gather*}
						And apply $Hom(X,-)$, surjectivity of $Hom(X,A)\rightarrow Hom(X,B)$.
					\end{itemize}
				\end{proof}
				\textbf{Note} Without the projective property, $Hom(X,-)$ is just left-exact. See discussion here:
				\subsubsection{Projective Resolution}
				\subsection{Injective Object}
					\begin{Def}
						An object is \textbf{injective} if it satisfies the right lifting property for \textbf{monomorphism}
						\begin{center}
							\begin{tikzcd}
								B\arrow[r]&A\\
								P\arrow[u,dashrightarrow]\arrow[ru,dashrightarrow]
							\end{tikzcd}
						\end{center}
					\end{Def}
					\subsubsection{Preservation of Injective Objects}
					\begin{lemma}
						Given a pair of additive adjoint functors
						\begin{gather*}
							(L)
						\end{gather*}
						betweeen abelian categories such that the left adjoint $L$ is a \textbf{left exact functor}, then the right adjoint preserve injective objects.
					\end{lemma}
					\begin{proof}
						内容...
					\end{proof}
					
					\begin{proposition}
						内容...
					\end{proposition}
					\begin{proof}
						内容...
					\end{proof}
					\subsubsection{Existence of Enough Injectives}
					Given some categories with enough injectives:
					\begin{itemize}
						\item 
						\item 
					\end{itemize}
					
					\begin{lemma}
						内容...
					\end{lemma}
					\begin{proof}
						内容...
					\end{proof}
					\subsubsection{Constructions}
					\begin{proposition}
						\begin{itemize}
							\item The (small) product of injective objects is injective;
							\item The direct sum of \underline{finite} injective objects is injective;
						\end{itemize}
					\end{proposition}
					\begin{proof}
						\begin{itemize}
							\item Consider the following diagram
							\begin{center}
								\begin{tikzcd}
									&\prod I_i\arrow[d]\\
									&I_i\arrow[d,twoheadrightarrow]\\
									A\arrow[r]\arrow[ru]\arrow[ruu,dashrightarrow]&B
								\end{tikzcd}
							\end{center}
							\item \begin{center}
								\begin{tikzcd}
									&\prod_{i\in I} I_i\arrow[dd,shift left=2,twoheadrightarrow,bend left=20]\\
									&\sum I_i\arrow[u]\arrow[d,two heads]\\
									A\arrow[r]\arrow[ru]&B
								\end{tikzcd}
							\end{center}
						\end{itemize}
					\end{proof}
					And the injectivity of arbitrary direct sum can be judged with Baer's criterion.
					\subsubsection{Injective Resolution}
					\begin{Def}
						An \textbf{injective} resolution is a cochain complex $J^*\in Ch^*(\mathcal{A})$ equipped with a \textbf{quasi-isomorphism}
						\begin{gather*}
							i:X\stackrel{\sim}{\rightarrow} J^*
						\end{gather*}
						such that $J^n\in \mathcal{A}$ is an injective object for all $n\in \NN$. i.e.Quasi-isomorphism
						\begin{center}
							\begin{tikzcd}
								X\arrow[r]\arrow[d]&0\arrow[r]\arrow[d]&\dots\arrow[r]&0\arrow[d]\arrow[r]&\dots\\
								J^0\arrow[r]&J^1\arrow[r]&\dots\arrow[r]&J^n\arrow[r]&\dots
							\end{tikzcd}
						\end{center}
					\end{Def}
					
					\begin{proposition}
						If $\mathcal{A}$ has enough injectives, then every $X\in \mathcal{A}$ has an injective resolution. 
					\end{proposition}
					\begin{proof}
						We have
					\end{proof}
					
					\begin{proposition}
						内容...
					\end{proposition}
					\begin{proof}
						内容...
					\end{proof}	\subsubsection{Injective Resolution}
					\begin{Def}
					For $X\in \mathcal{A}$ an object, an \textbf{injective resolution} is a cochain complex $J^*\in Ch^*(\mathcal{A})$ equipped with a quasi-isomorphism
					\begin{gather*}
						i:X\stackrel{\sim}{\rightarrow} J^*
					\end{gather*}
					such that $J^n\in \mathcal{A}$ is an \textbf{injective object} for all $n\in \NN$.
					\end{Def}
					\begin{center}
					\begin{tikzcd}
						X\arrow[r]\arrow[d,"i^0"]&0\arrow[r]\arrow[d]&\dots\arrow[r]&0\arrow[d]\arrow[r]&\dots\\
						J^0\arrow[r,"d^0"]&J^1\arrow[r,"d^1"]&\dots\arrow[r]&J^n\arrow[r,"d^n"]&\dots
					\end{tikzcd}
					\end{center}
					\subsubsection{Projective Resolution}
					\begin{Def}
					For $X\in \mathcal{A}$ is an object, a \textbf{projective resolution} of $X$ is a chain ccomplex $J_*\in Ch_\bullet(\mathcal{A})$ equipped with a quasi-isomorphism
					\begin{gather*}
						p:J_*\stackrel{\sim}{\rightarrow} X
					\end{gather*}
					such that $J_n\in \mathcal{A}$ is a projective object for all $n\in \NN$.
					\end{Def}
					\begin{center}
					\begin{tikzcd}
						\dots\arrow[r]\arrow[d]&J_n\arrow[r]\arrow[d]&\dots\arrow[r]&J_1\arrow[r]\arrow[r]\arrow[d]&J_0\arrow[d]\\
						\dots\arrow[r]&0\arrow[r]&\dots\arrow[r]&0\arrow[r]&X
					\end{tikzcd}
					\end{center}
					
					\begin{proposition}
					Let $\mathcal{A}$ be an abelian category with enough injectives, such as $R-Mod$ for some ring $R$, then every object $X\in \mathcal{A}$ has an injective resolution.
					\end{proposition}
					\begin{proof}
					内容...
					\end{proof}
					\begin{proposition}
					Let $\mathcal{A}$ be an abelian category with enough projectives(such as $R$-Mod for some ring $R$). Then every object $X\in \mathcal{A}$ has a projective resolution.
					\end{proposition}
					\begin{proof}
					内容...
					\end{proof}
					\subsection{Resolutions, and $Ext^1$}
					Suppose $\mathcal{A}$ be an abelian category.
					\subsubsection{Left Resolution}
					
					\begin{Def}
					A \textbf{left resolution} of $M\in \mathcal{A}$ consists of:
					\begin{itemize}
						\item Chain complex $A_\bullet$;
						\item Morphism $\epsilon:A_0\rightarrow M$;
						\item Exactness
						\begin{gather*}
							\dots\rightarrow A_2\rightarrow A_1\rightarrow A_0\stackrel{\epsilon}{\rightarrow} M\rightarrow 0
						\end{gather*}
					\end{itemize}
					\end{Def}
					\subsubsection{Free Resolution}
					\begin{Def}
					Let $M\in \mathcal{A}$. A \textbf{free resolution}
					\end{Def}
					Conclusions
					\begin{proposition}
					Every finitely generated abelian group admits a free resolution with length 2.
					\end{proposition}
					\begin{proof}
					Take the free object generated by
					\end{proof}
					Some of contents will be given in:Syzygies.
			\subsection{Universally Injective Module Maps}
				\begin{Def}
					\begin{itemize}
						\item Let $f:M\rightarrow N$ be a map of $R$-modules.
						
						Then $f$ is called \textbf{universally injective} if for every $R$-module $Q$, the map
						\begin{gather*}
							f\otimes_R id_Q:M\otimes_R Q\rightarrow N\otimes_R Q
						\end{gather*}
						
						is injective.
						\item A sequence $0\rightarrow M_1\rightarrow M_2\rightarrow M_3\rightarrow 0$ of $R$-modules is called \textbf{universally exact} if it is exact, and $M_1\rightarrow M_2$ is universally injective.
					\end{itemize}
				\end{Def}
				
				\begin{example}
					内容...
				\end{example}
				
				\begin{theorem}
					Let 
					\begin{gather*}
						0\rightarrow M_1\stackrel{f_1}{\rightarrow} M_2 \stackrel{f_2}{\rightarrow} M_3\rightarrow0
					\end{gather*}
					be an exact sequence of $R$-modules. The following are equivalent:
					\begin{itemize}
						\item Sequence $0\rightarrow M_1\rightarrow M_2\rightarrow M_3\rightarrow 0$ is universally exact;
						\item For every finitely presented $R$-module $Q$, the sequence
						\begin{gather*}
							0\rightarrow M_1\otimes_R Q\rightarrow M_2\otimes_R Q\rightarrow M_3\otimes_R Q\rightarrow 0
						\end{gather*}
						\item Given elements $x_i\in M_1;i=1,\dots,n;y_i\in M_2,j=1,\dots,m$. And $a_{ij}\in R(i=1,\dots,n,j=1,\dots,m)$ such that for all $i$
						\begin{gather*}
							f_1(x_i)=\sum_{j} a_{ij}y_j
						\end{gather*}
						
						there exists $z_i\in M_1$ such that for all $i$:
						\begin{gather*}
							内容...
						\end{gather*}
						\item Given a commutative diagram of $R$-module maps
						\begin{center}
							\begin{tikzcd}
								内容...
							\end{tikzcd}
						\end{center}
						\item 
						\item 
					\end{itemize}
				\end{theorem}
				\begin{proof}
					内容...
				\end{proof}
				
				\begin{lemma}
					Let
					\begin{gather*}
						0\rightarrow M_1\rightarrow M_2\rightarrow M_3\rightarrow 0
					\end{gather*}
					
					be an exact sequence of $R$-modules. Suppose that: $M_3$ is \underline{finite presentaion}. Then, 
					\begin{gather*}
						内容...
					\end{gather*}
				\end{lemma}
				\begin{proof}
					内容...
				\end{proof}
		\section{Flatness}
			\subsection{Exact Sequence}
				
				\begin{proposition}
					(Exactness of Flatness) Suppose $0\rightarrow A\rightarrow B\rightarrow C\rightarrow 0$ an exact sequence with $A,C$ flat. Then, $B$ is flat.
				\end{proposition}
				\begin{proof}
					We just consider the following diagram: Notice that the tensor product $-\otimes X$ is in general right-exact, $X=A,B,C$ we just need to draw the items, for the short exact sequence $0\rightarrow M'\rightarrow M$ we have
					\begin{center}
						\begin{tikzcd}
							&M'\otimes A\arrow[r]\arrow[d,"f"]&M'\otimes B\arrow[r]\arrow[d,"g"]&M'\otimes C\arrow[r]\arrow[d,"h"]&0\\
							0\arrow[r]&M\otimes A\arrow[r]&M\otimes B\arrow[r]&M\otimes C
						\end{tikzcd}
					\end{center}
					
					We can get an exact sequence $ker(f)\rightarrow ker(g)\rightarrow ker(h)$, and $ker(f)=ker(h)=0\Rightarrow ker(g)=0$.
					
					Another solution is to use the statement with 9-lemma:
					\begin{center}
						\begin{tikzcd}
							&0\arrow[d]&0\arrow[d]&0\arrow[d]\\
							0\arrow[r]&M'\otimes A\arrow[r]\arrow[d]&M'\otimes B\arrow[r]\arrow[d]&M'\otimes C\arrow[r]\arrow[d]&0\\
							0\arrow[r]&M\otimes A\arrow[r]\arrow[d]&M\otimes B\arrow[r]\arrow[d]&M\otimes C\arrow[r]\arrow[d]&0\\
							0\arrow[r]&M''\otimes A\arrow[r]\arrow[d]&M''\otimes B\arrow[r]\arrow[d]&M''\otimes C\arrow[r]\arrow[d]&0\\
							&0&0&0
						\end{tikzcd}
						
						Column 1,3 are exact, so is column 2.
					\end{center}
				\end{proof}
			\subsection{Characterizing Flatness}
				\begin{lemma}
					Let $M$ be an $R$-module. The following are equivalent:
					\begin{itemize}
						\item $M$ is flat;
						\item If $f:R^n\rightarrow M$ is a module map, $x\in Ker(f)$, then there are module maps
						\begin{gather*}
							h:R^n\rightarrow R^m\qquad g:R^m\rightarrow M
						\end{gather*}
						such that $f=g\circ h$ and $x\in Ker(h)$.
						\item 
						\item 
					\end{itemize}
				\end{lemma}
				\begin{proof}
				
				\end{proof}
				
				\begin{theorem}
					(Lazard's Theorem) Let $M$ be an $R$-module. Then,
					
					$M$ is flat$\Leftrightarrow$ It is the colimit of a directed system of free finite $R$-modules.
				\end{theorem}
				\begin{proof}
					内容...
				\end{proof}
				
				Moreover, 
				\begin{itemize}
					\item https://arxiv.org/pdf/1109.0439 'The Lazard-Linke's Theorem on quasi-coherent sheaves.
				\end{itemize}
			\subsection{Flatness over PID}
				
	\chapter{Chain Complexes}
		\section{Some of Category Theory}
			\subsection{$\mathbf{Ab}$-Category}
			\begin{Def}
				A category $\mathcal{A}$ is called an $\mathcal Ab$-category if every $hom$-set $Hom_\mathcal{A}(A,B)$ in $\mathcal{A}$ is given the structure of an abelian group in such a way that composition distributes over addition.
				
				In particular, given a diagram in $\mathcal{A}$ of the form
				\begin{center}
					\begin{tikzcd}
						A\arrow[r,"f"]&B\arrow[r,shift left=3,"g'"]\arrow[r,shift right=3,"g",swap]&C\arrow[r,"h"]&D
					\end{tikzcd}
				\end{center}
				We have
				\begin{gather*}
					h(g+g')f=hgf+hg'f
				\end{gather*}
			\end{Def}
			\subsection{Preadditive Category}
				(Stack Project, Homological Algebra, 3)
				\begin{Def}
					A category $\mathcal{A}$ is \textbf{preadditive} if each morphism set $Mor_\mathcal{A}(x,y)$ is endowed with a structure of an abelian group such that:
					\begin{itemize}
						\item  Each composition
						\begin{gather*}
							Mor(x,y)\times Mor(y,z)\rightarrow Mor(x,z)
						\end{gather*}
						is bilinaer($(f+g)\circ h=f\circ h+g\circ h,f\circ (g+h)=f\circ g+f\circ h$)
					\end{itemize}
				\end{Def}
				Additive functor
				\begin{Def}
					\begin{itemize}
						\item An addive functor $F:\mathcal{B}\rightarrow \mathcal{A}$ between $\mathcal Ab$-categories $\mathcal{B}$ and $\mathcal{A}$ is a functor such that each
						\begin{gather*}
							Hom_\mathcal{B}(B',B)\rightarrow Hom_{\mathcal{A}}(A',A)
						\end{gather*}
					\end{itemize}
					is a group homomorphism;
				\end{Def}
				\begin{lemma}
					Let $\mathcal{A}$ be a preadditive. Let $x\in Ob(\mathcal{A})$, the following are equivalent:
					\begin{itemize}
						\item $x$ is an initial object;
						\item $x$ is a final object;
						\item $id_x=0\in Mor_\mathcal{A}(x,x)$
					\end{itemize}
					Furthermore, if such $0$ exists, then: A morphism $\alpha:y\rightarrow z$ factors through $0\Leftrightarrow \alpha=0$
				\end{lemma}
				\begin{proof}
					内容...
				\end{proof}
				Such object is called \textbf{zero object} in preadditive category.
				
				\textbf{Remark} Existence:
				\begin{lemma}
					(Isomorphism of Finite Product, Coproduct) In preadditive category 
					\begin{itemize}
						\item $x\times y$ exists$\Leftrightarrow x\sqcup y$ exists;
						\item 
						 we have
						\begin{gather*}
							x\sqcup y\cong x\times y
						\end{gather*}
					\end{itemize}
				\end{lemma}
				\begin{proof}
					\begin{itemize}
						\item If $z=x\sqcup y$, define $p:=i\sqcup 0,p:=i'\sqcup 0$
						\item If $z=x\times y$, define $i:=id_x\times 0,i'=0\times id_y$
					\end{itemize}
					\begin{center}
						\begin{tikzcd}
							x\arrow[rr,"id_x"]\arrow[rd,"i"]&&x\\
							&z\arrow[rd,"p'"]\arrow[ru,"p"]\\
							y\arrow[rr,"id_y"]\arrow[ru,"i'"]&&y
						\end{tikzcd}
					\end{center}
				\end{proof}
				\begin{Def}
					(Direct Sum) In preadditive category, such $x\oplus y:=x\times y=x\sqcup y$.
				\end{Def}
				
				\begin{lemma}
					
				\end{lemma}
				\begin{proof}
					内容...
				\end{proof}
			\subsection{Additive Funtors, Additive Category}
			
			\begin{Def}
				\begin{itemize}
					\item An \textbf{additive category} is an $\mathcal Ab$-category $\mathcal{A}$ with:
					\begin{itemize}
						\item Zero Object $0\in \mathcal{A}$;
						\item Product $A\times B$;
					\end{itemize}
				\end{itemize}
			\end{Def}
			\begin{lemma}
				For additive category, we have
				\begin{gather*}
					A\times 0\cong A, A\in \mathcal{A}
				\end{gather*}
			\end{lemma}
			\begin{proof}
				Consider 
				\begin{center}
					\begin{tikzcd}
						&A\\
						A\times 0\arrow[ru,"p_A"]\arrow[rd]&&A\arrow[ll,"f"]\arrow[lu,"id_A"]\arrow[ld]\\
						&0
					\end{tikzcd}
				\end{center}
				We have $p_A\circ f=id_A,f\circ p_A=id_{A\times 0}$
			\end{proof}
			\begin{lemma}
				(Product Map) If $f\in Mor_\mathcal{A}(B,C),g\in Mor_\mathcal{A}(D,E)$, then we have
				\begin{gather*}
					f\times g:B\times D\rightarrow C\times E
				\end{gather*} 
			\end{lemma}
			\begin{proof}
				Consider following diagram
				\begin{center}
					\begin{tikzcd}
						B\arrow[r]&C\\
						B\times D\arrow[r,dashrightarrow]\arrow[u]\arrow[d]&C\times E\arrow[u]\arrow[d]\\
						D\arrow[r]&E
					\end{tikzcd}
				\end{center}
			\end{proof}
			\begin{proposition}
				If $\mathcal{A}$ is an additive category, then the product exists, with
				\begin{gather*}
					A\times B=A\sqcup B
				\end{gather*}
			\end{proposition}
			\begin{proof}
				Firstly, $A\times B$ has following mappings $\iota_A:A\cong A\times 0\rightarrow A\times B$(B resp.); And then for $p_A:P\rightarrow A,p_B:P\rightarrow B$ we have
				\begin{gather*}
					\iota_A p_A+\iota_Bp_B:P\rightarrow A\times B
				\end{gather*}
				Thus, $A\times B$ is the coproduct.
			\end{proof}
			Also take the following notion
			\begin{gather*}
				A\times B\cong A\oplus B
			\end{gather*}\label{Homological-Algebra-Abelian-Category-Product-Equals-Coproduct}
			
			Then,
			\begin{proposition}
				The following statements are equivalent:
				\begin{itemize}
					\item Category $\mathcal{C}$ is additive;
					\item The structure maps
					\begin{itemize}
						\item \begin{gather*}
							i_X:X\rightarrow X\oplus Y\\
							i_Y:Y\rightarrow X\oplus Y
						\end{gather*}
						\item \begin{gather*}
							p_X:X\oplus Y\rightarrow X\\
							p_Y:X\oplus Y\rightarrow Y
						\end{gather*}
					\end{itemize}
					with
					\begin{itemize}
						\item \begin{gather*}
							p_X\circ i_X=1_X\qquad p_Y\circ i_Y=1_Y\\
							p_X\circ i_Y=0\qquad p_Y\circ i_X=0
						\end{gather*}
						\item Completeness
						\begin{gather*}
							1_{X\oplus Y}=i_X\circ p_X+i_Y\circ p_Y
						\end{gather*}
					\end{itemize}
				\end{itemize}
			\end{proposition}
			\begin{proof}
				Omitted.
			\end{proof}
			
			In additive category, the morphism between direct products can be represented as matrix
			\begin{gather*}
				f:A\oplus B\rightarrow C\oplus D\\
				f=\begin{pmatrix}
					p_W\circ f\circ i_X&p_W\circ f\circ i_Y\\
					p_Z\circ f\circ i_X&p_Z\circ f\circ i_Y
				\end{pmatrix}=\begin{pmatrix}
					f_{11}&f_{12}\\
					f_{21}&f_{22}
				\end{pmatrix}
			\end{gather*}
			\subsection{Karoubian Categories}
			(Stack Project, HA, 4) 
			
			\begin{Def}
				Let $\mathcal{C}$ be a preadditive category. 
				
				We say $\mathcal{C}$ is \textbf{Karoubian} if every idempotent endomorphism $(f^2=f)$ of an object of $\mathcal{C}$ has a kernel.
				\begin{gather*}
					f^2=f\Rightarrow \exists ker(f),(0\rightarrow ker(f)\rightarrow x\rightarrow y)
				\end{gather*}
			\end{Def}
			\begin{proposition}
				(Equivalent Conditions for Karoubian Categories) The followings are equivalent:
				\begin{itemize}
					\item $\mathcal{C}$ is Karobian;
					\item Each idempotent endomorphism of an object $\mathcal{C}$ has a cokernel;
					\item Given an idempotent endomorphism $p:z\rightarrow z$, there exists a direct sum decomposition $z=x\oplus y$ such that: $p$ corresponds to the projecion onto $y$
				\end{itemize}
			\end{proposition}
			\begin{proof}
				\begin{itemize}
					\item $(1)\Rightarrow (2)$
					\item $(2)\Rightarrow (3)$
					\item $(3)\Rightarrow (1)$
				\end{itemize}
			\end{proof}
			
			\begin{proposition}
				(Criterias for Karoubian Category) Let $\mathcal{D}$
				\begin{itemize}
					\item If $\mathcal{D}$ has countable products and kernels of maps which have a right inverse, then Karoubian; $\mathcal{D}$
					\item If $\mathcal{D}$ has countable ciproduct and cokernels of maps which have a left inverse, then Karoubian; $\mathcal{D}$
				\end{itemize}
			\end{proposition}
			\begin{proof}
				\begin{itemize}
					\item 
					\item 
				\end{itemize}
			\end{proof}
			\textbf{Warning} Preadditive Category $\subset$ Karoubian Category$\subset$ Additive Category
			
			\begin{example}
				Category of free modules.
			\end{example}
			
			\begin{example}
				Category $\mathcal Set$.
			\end{example}
			\subsection{Abelian Category}
			\begin{Def}
				Define: For $f\in Mor_\mathcal{A}(B,C)$,
				\begin{itemize}
					\item (Kernel) The \textbf{kernel} is the universal object given by $i:Ker(f)\rightarrow B$ with $fi=0$
					\item (Cokernel) The \textbf{cokernel} is the universal object given by $j:C\rightarrow Coker(f)$ with $jf=0$.
				\end{itemize}
				And:
				\begin{itemize}
					\item $f$ is mono iff $fg=fh\Rightarrow g=h$
					\item $f$ is epi iff $gf=hf\Rightarrow g=h$
				\end{itemize}
			\end{Def}
			Then,
			\begin{Def}
				The \textbf{abelian category} satisfies:
				\begin{itemize}
					\item $f$ is  Mono/injective$\Leftrightarrow$ $Ker(f)=0$;
					\item $f$ is Epi/surjecitve$\Leftrightarrow Coker(f)=0$;
				\end{itemize}
			\end{Def}
			Define:
			\begin{gather*}
				Coim(f):=Coker(Ker(f)\rightarrow B)\\
				Im(f):=Ker(C\rightarrow Coker(f))
			\end{gather*}
			
			\textbf{Remark} Sometimes take the notion:
			\begin{itemize}
				\item inj.:$Ker(f)=0$;
				\item sur.:$Coker(f)=0$;
			\end{itemize}
			
			\begin{proposition}
				Category $\mathcal{A}$ is abelian iff it is additive,  all kernels and cokernels exist, and the natural map
				\begin{gather*}
					Coim(f)\rightarrow Im(f)
				\end{gather*}
				is an isomorphism for all morphisms $f$ of $\mathcal{A}$.
			\end{proposition}
			\begin{proof}
				Notice: $Coim(f)=Coker(Ker(f)\rightarrow B)=Ker(C\rightarrow Coker(f))=Im(f)$, now consider following property
				\begin{gather*}
					Ker(f)\rightarrow B\rightarrow Coim(f)=Im(f)\rightarrow C\rightarrow Coker(f)
				\end{gather*}
				Then, equivalent by universal property. Omitted.
			\end{proof}
			
			We could define the exactness:
			\begin{Def}
				For abelian category $\mathcal{A}$, we could define the \textbf{exactness}
				\begin{gather*}
					0\rightarrow A\stackrel{f}{\rightarrow} B\stackrel{g}{\rightarrow} C\rightarrow 0
				\end{gather*}
				With $Im(f)=Ker(g)$.
			\end{Def}
			
			\begin{lemma}
				Let: $\mathcal{A}$ be a preadditive category.
				
				The addition on sets of morphisms make $\mathcal{A}^{opp}$ into a preadditive category.
				
				Furthermore, $\mathcal{A}$ is preadditive/additive/abelian$\Leftrightarrow \mathcal{A}^{op}$ is preadditive/additive/abelian	;
			\end{lemma}
			\begin{proof}
				内容...
			\end{proof}
			
			\begin{lemma}
				(Finite Limimts, Colimits) The finite limits, colimits exist in the abelian category.
			\end{lemma}
			\begin{proof}
				\begin{itemize}
					\item Limit
					\item Colimit
				\end{itemize}
			\end{proof}
			
			\begin{theorem}
				内容...
			\end{theorem}
			\begin{proof}
				内容...
			\end{proof}
			\subsection{Pullback,Pushout Diagram}
					\begin{Def}
						The \textbf{intersection/pullback} of objects $x,y\in \mathcal{A}$, which satisfies the following diagram
						\begin{center}\begin{tikzcd}
								x \cap y \arrow[r, "p_1"] \arrow[d, "p_2"] \arrow[rd, phantom, "\lrcorner", very near start] & x \arrow[d, hookrightarrow] \\
								y \arrow[r, hookrightarrow] & z
							\end{tikzcd}
						\end{center}
						And the \textbf{cointersection/pushout}
						\begin{center}
							\begin{tikzcd}
								z \arrow[r,twoheadrightarrow] \arrow[d,twoheadrightarrow] & x \arrow[d] \\
								y \arrow[r] & x \cup_z y
							\end{tikzcd}
						\end{center}
					\end{Def}
				\subsection{Lemmas}
					\subsubsection{3-Lemma}
					\subsubsection{Snake Lemma}
						\begin{proposition}
							内容...
						\end{proposition}
						\begin{proof}
							内容...
						\end{proof}
					\subsubsection{4-Lemma}
					\subsubsection{5-Lemma}
				\subsection{Grothendieck Category}
				
				\url{https://ncatlab.org/nlab/show/additive+and+abelian+categories}
				
				\url{https://ncatlab.org/nlab/show/Grothendieck+category}
				\begin{Def}
					(Grothendieck) 
					\begin{itemize}
						\item (AB1) Existence of kernel,cokernel;
						\item (AB2) Isomorphism $Coim(f)\rightarrow Im(f)$,where $Coim(f)=Coker(Ker(f)\rightarrow A),Im(f)=Ker(A\rightarrow Coker(f))$;
						
						(AB1)+(AB2)=Abelian
						\item (AB3) An abelian category with all coproducts; Hence all colimits;
						\item (AB4) AB3 Satisfied, coproduct of monomorphism is monic;
						\item (AB5) AB3 Satisfied, filtered limits of exact seuqneces \underline{exist} and are \underline{exact};
						\item (AB6) AB3 Satisfied, such that for increasing directed set $J$, $(B^j)_{j\in J}$ and families $B^j=(B^j_i)_{i\in I_j}$(Thus, $B=((B^j_i)_{i\in I_j})_{j\in J}$) then, we have
						\begin{gather*}
							\cap_{j\in J} (\sum_{i\in I_j} B^i_j)=\sum_{(i_j)\in \prod I_j} (\cap_{j\in J} B^i_{i_j})
						\end{gather*}
						\item (AB$3^*$) DF
						\item (AB$4^*$) DF
						\item (AB$5^*$) DF
						\item (AB$6^*$) DF	
					\end{itemize}
				\end{Def}
				\subsection{Extension}
					\begin{Def}
						\begin{itemize}
							\item An \textbf{extension} of $A$ by $B$ in abelian category $\mathcal{A}$ is an (short) exact sequence
							\begin{gather*}
								0\rightarrow A\rightarrow B\rightarrow C\rightarrow 0
							\end{gather*}
							\item The morphism between extensions: For short exact sequences $0\rightarrow A\rightarrow E\rightarrow B\rightarrow 0;0\rightarrow A\rightarrow F\rightarrow B\rightarrow 0$. And the following diagram commutes:
							\begin{center}
								\begin{tikzcd}
									0&A&E&B&0\\
									0&A&F&B&0
								\end{tikzcd}
							\end{center}
							\item The isomorphic class:Let $A,B\in Ob(\mathcal{A})$, the set of isomorphic classes of extension of $B$ by $A$ denoted
							\begin{gather*}
								Ext_\mathcal{A}(B,A)
							\end{gather*}
							is called the \textbf{Ext-group}.
						\end{itemize}
					\end{Def}
					
					Actually, we turn $Ext(-,-)$ into functor
					\begin{itemize}
						\item Pullback
						\item Pushout
					\end{itemize}
					
					
					
					\begin{lemma}
						The construction
						\begin{gather*}
							(E_1,E_2)\mapsto E_1+E_2
						\end{gather*}
						defines a commutative group law on $Ext_\mathcal{A}(B,A)$, which is functorial on each variable.
					\end{lemma}
					\begin{proof}
						内容...
					\end{proof}
					
					\begin{proposition}
						(Six-Item Exact Sequence) There exists long exact sequence
						\begin{center}
							\begin{tikzcd}
								0&Hom(C,X)&Hom(B,X)&Hom(A,X)\\
								&Ext(C,X)&Ext(B,X)&Ext(A,X)&0
							\end{tikzcd}
						\end{center} 
					\end{proposition}
					\begin{proof}
						内容...
					\end{proof}
					
					\begin{proposition}
						(Six-Item Exact Seuquence)
						\begin{center}
							\begin{tikzcd}
								0&Hom(X,A)&Hom(X,B)&Hom(X,C)\\
								&Ext(X,A)&Ext(X,B)&Ext(X,C)&0
							\end{tikzcd}
						\end{center}
					\end{proposition}
					\begin{proof}
						内容...
					\end{proof}
				\subsection{Additive Functors}
					\begin{lemma}
						Let $\mathcal{A},\mathcal{B}$ \underline{abelian categories}, $F:\mathcal{A}\rightarrow \mathcal{B}$ is functor. Then, the followings are equivalent
						\begin{itemize}
							\item 
							\item 
							\item 
						\end{itemize}
					\end{lemma}
					\begin{proof}
						内容...
					\end{proof}
					
					\begin{lemma}
						内容...
					\end{lemma}
					\begin{proof}
						内容...
					\end{proof}
				\subsection{Localization}
				\subsection{Jordan-Hölder}
				\subsection{Serre Subcategories}
					\begin{Def}
						Let $\mathcal{A}$ be an abelian category:
						\begin{itemize}
							\item A \textbf{Serre subcategory} of $\mathcal{A}$ is a \underline{nonempty full subcategory} $\mathcal{C}$ such that for every exact sequence
							\begin{center}
								\begin{tikzcd}
									A\arrow[r]&B\arrow[r]&C
								\end{tikzcd}
							\end{center}
							with $A.C\in Ob(\mathcal{A})$ such that
							\begin{gather*}
								A,C\in Ob(\mathcal{C})\Rightarrow B\in Ob(\mathcal{C})
							\end{gather*}
							\item A \textbf{weak Serre subcategory} $\mathcal{A}$ is a nonempty full subcategory $\mathcal{C}$ of $\mathcal{A}$ such that: Given any exact sequence
							\begin{gather*}
								A_0\rightarrow A_1\rightarrow A_2\rightarrow A_3\rightarrow A_4
							\end{gather*}
							with $A_0,A_1,A_3,A_4\in Ob(\mathcal{C})\Rightarrow A_2\in Ob(\mathcal{C})$.
						\end{itemize}
					\end{Def}
					
					\begin{lemma}
						(Equivalent Condition of Serre Subcategory) Let $\mathcal{A}$ be an abelian category, $\mathcal{C}$ be a subcategory of $\mathcal{C}$. Then, the following condition are equivalent: $\mathcal{C}$ is Serre subcategory$\Leftrightarrow$ \begin{itemize}
							\item $0\in Ob(\mathcal{C})$;
							\item $\mathcal{C}$ is a strictly full subcategory of $\mathcal{A}$;
							\item Any subobject or quotient of $\mathcal{C}$ is an object of $\mathcal{C}$;
							\item If $A\in Ob(\mathcal{A})$ is an extension of objects of $\mathcal{C}$, then $A\in Ob(\mathcal{C})$.
						\end{itemize}
					\end{lemma}
					\begin{proof}
						$\Leftarrow$ 
						
						$\Rightarrow$ Omitted.
					\end{proof}
					
					\begin{lemma}
						(Equivalent Condition of Weak Serre Subcategories)
						\begin{itemize}
							\item 
							\item 
							\item 
							\item 
						\end{itemize}
					\end{lemma}
					\begin{proof}
						内容...
					\end{proof}
					\subsubsection{Serre Subcategory by Morphism}
						\begin{lemma}
							Let $\mathcal{A},\mathcal{B}$ be abelian categories, and $F:\mathcal{A}\rightarrow \mathcal{B}$ be an exact functor.
							
							Then, the \textbf{full subcategory} specified by $F(\mathcal{C})=0$ forms a Serre subcategory of $A$.
						\end{lemma}
						\begin{proof}
							Notice that the functor is exact.
						\end{proof}
						
						\begin{Def}
							Let $\mathcal{A},\mathcal{B}$ be an exact functors, then the full subcategory of objects $\mathcal{C}$ such that $F(\mathcal{C})=0$ is called the \textbf{kernel} of the functor $F$. Sometimes denoted by
							\begin{gather*}
								ker(\mathcal{F})
							\end{gather*} 
						\end{Def}
					\subsubsection{Quotient}
						\begin{lemma}
							(Categoricial Quotient) Let $\mathcal{A}$ be an abelian category, $\mathcal{C}\subseteq\mathcal{A}$ be a Serre subcategory, there exists an abelian category and an exact functor
							\begin{gather*}
								F:\mathcal{A}\rightarrow \mathcal{A}/\mathcal{C}
							\end{gather*}
							which is \textbf{essentially surjective} and whose kernel is $\mathcal{C}$.
							
							The category $\mathcal{A}/\mathcal{C}$ is specified by the following universal property: For any exact functor
							\begin{gather*}
								G:\mathcal{A}\rightarrow \mathcal{B}
							\end{gather*} 
							Such that $\mathcal{C}\subseteq Ker(\mathcal{G})$, there exists a factorization $G=H\circ F$ for a unique exact functor $H:\mathcal{A}/\mathcal{C}\rightarrow \mathcal{B}$.
						\end{lemma}
						\begin{proof}
							内容...
						\end{proof}
						
						\begin{lemma}
							Let $\mathcal{A},\mathcal{B}$ be abelian categories, $F:\mathcal{A}\rightarrow \mathcal{B}$ be an exact functor, let $\mathcal{C}\subseteq \mathcal{A}$ be Serre subcategory contained in the kernel of $F$. Then,
							\begin{gather*}
								\mathcal{C}=Ker(F)\Leftrightarrow \bar{F}:\mathcal{A}/\mathcal{C}\rightarrow \mathcal{B}\mbox{ is faithful}
							\end{gather*}
						\end{lemma}
						\begin{proof}
							内容...
						\end{proof}
				\subsection{K-Groups of Abelian Category}
					\subsubsection{0th Abelian Category}
					\begin{Def}
						Let $\mathcal{A}$ be an abelian category, then, denote $K_0(\mathcal{A})$ the \textbf{zeroth K-group} of $\mathcal{A}$ and for exact sequence
					\end{Def}
					\textbf{Remark} Another way, we to say this is that: There exists a presentation
					\begin{gather*}
						\oplus_{(A\rightarrow B\rightarrow C)\in Ses} Z[A\rightarrow B\rightarrow C]\rightarrow \oplus_{A\in Ob(\mathcal{A})} Z[A]\rightarrow K_0(A)\rightarrow 0
					\end{gather*}
					\begin{lemma}
					\end{lemma}
					\begin{proof}
						内容...
					\end{proof}
					
					\begin{lemma}
						内容...
					\end{lemma}
					\begin{proof}
						内容...
					\end{proof}
					
					\begin{lemma}
						内容...
					\end{lemma}
					\begin{proof}
						内容...
					\end{proof}
				\subsection{Cohomological $\delta$-Functor}
					\begin{Def}
						Let $\mathcal{A},\mathcal{B}$ be abelian categories, a \textbf{cohomological $\delta$-functor} or simly \textbf{$\delta$-functor} from $\mathcal{A}\rightarrow \mathcal{B}$ is given by the following datas:
						\begin{itemize}
							\item A collection 
							\begin{gather*}
								F^n:\mathcal{A}\rightarrow \mathcal{B},n\geq 0
							\end{gather*}
							be additive functors, and
							\item For any short exact sequence $0\rightarrow A\rightarrow B\rightarrow C\rightarrow 0$ of $\mathcal{A}$, a collection 
							\begin{gather*}
								\delta_{A\rightarrow B\rightarrow C}:F^n(C)\rightarrow F^{n+1}(A),n\geq 0
							\end{gather*}
							of morphisms of $\mathcal{B}$.
						\end{itemize}
						These data are assumed to satisfy the following axioms:
						\begin{itemize}
							\item For any short exact sequence as above, we have the \textbf{long exact sequence}
							\begin{center}
								\begin{tikzcd}
									0\arrow[r]&F^0(A)\arrow[r]&F^0(B)\arrow[r]&F^0(C)\arrow[lld]\\
									&F^1(A)\arrow[r]&F^1(B)\arrow[r]&F^1(C)\arrow[lld]\\
									&\dots
								\end{tikzcd}
							\end{center}
							\item For every morphism of short exact sequence $(0\rightarrow A\rightarrow B\rightarrow C\rightarrow 0)\rightarrow (0\rightarrow A'\rightarrow B'\rightarrow C'\rightarrow 0)$ of short exact sequences, the diagram
							\begin{center}
								\begin{tikzcd}
									F^n(C)\arrow[r,"\delta_{A\rightarrow B\rightarrow  C}"]\arrow[d]& F^{n+1}(A)\arrow[d]\\
									F^n(C')\arrow[r,"\delta_{A'\rightarrow B'\rightarrow C'}"]&F^{n+1}(A')
								\end{tikzcd}
							\end{center}
						\end{itemize}
						This implies, that $F^0(-)$ is left-exact.
					\end{Def}
					\subsubsection{Morphism of $\delta$-Functors}
					\begin{Def}
						内容...
					\end{Def}
					\subsubsection{Universal $\delta$-Functor}
					\begin{Def}
						内容...
					\end{Def}
					
					\begin{lemma}
						内容...
					\end{lemma}
					\begin{proof}
						内容...
					\end{proof}
					
					\begin{lemma}
						内容...
					\end{lemma}
					\begin{proof}
						内容...
					\end{proof}
			\section{Chain Complex}
				\begin{Def}
					\begin{itemize}
						\item A $\ZZ$-\textbf{graded chain complex} in $R-Mod$ is:
						\begin{itemize}
							\item A collection of obejcts $\{C_n\}_{n\in \ZZ}$;
							\item Morphisms $\partial_n:C_n\rightarrow C_{n-1}$ such that
							\begin{gather*}
								\partial_n\circ \partial_{n+1}=0,n\in \NN
							\end{gather*}
							Then:
							\begin{itemize}
								\item $\partial_n$ are called the \textbf{differentials} or \textbf{boundary maps};
								\item The elements of $C_n$ are called $n$-chains;
								\item For $n\geq 1$ the elements in the kernel
								\begin{gather*}
									Z_n:=Z_n(C_\bullet)=Ker(\partial_{n})
								\end{gather*}
								are called \textbf{n-cycles};
								\item The image
								\begin{gather*}
									B_n:=B_n(C_\bullet)=im(\partial_{n+1})
								\end{gather*}
								of $\partial_n:C_{n}\rightarrow C_{n-1}$ are called the \textbf{n-boundaries};
							\end{itemize}
							There exists an inclusion $\partial^2=0\Rightarrow 0\hookrightarrow B_n\hookrightarrow Z_n\hookrightarrow C_n$; Define $H_n:=Z_n/B_n$, then the \textbf{chain homology} of $C_*$
							\begin{gather*}
								0\rightarrow B_n\rightarrow Z_n\rightarrow H_n\rightarrow 0
							\end{gather*}
						\end{itemize}
					\end{itemize}
				\end{Def}
				
				Then,
				\begin{Def}
					\begin{itemize}
						\item A \textbf{chain map} $f:V_*\rightarrow W_*$ is a collection of morphism $(f_n:V_n\rightarrow W_n)_{n\in \ZZ}$ in $\mathcal{A}$ such that all the diagram
						\begin{center}
							\begin{tikzcd}
								V_{n+1}\arrow[r]\arrow[d]&V_n\arrow[d]\\
								W_{n+1}\arrow[r]&W_n
							\end{tikzcd}
						\end{center}
						\item Take the notation of $Ch(R-Mod)$ for the category. Furthermore, for the category $\mathcal{A}$, take the notion of $Ch(\mathcal{A})$.
					\end{itemize}
				\end{Def}
				Especially, when $R=\ZZ$, is the chain complex of \underline{abelian groups}.
				\begin{proposition}
					$Ch(\mathcal{A})$ is an abelian category, iff $\mathcal{A}$ is an abelian category.
				\end{proposition}
				\begin{proof}
					$\Rightarrow$ Omitted. There exists a faithful funcotor $\mathcal{A}\rightarrow Ch(\mathcal{A})$;
					
					$\Leftarrow$ $\mathcal{A}$ is additive, with structure induced by $\mathcal{A}$. And the biproduct is well-defined. The kernel commutes with 
				\end{proof}
				
				\begin{corollary}
					$Ch(R-Mod)$ is an abelian category.
				\end{corollary}
				\begin{proof}
					Omitted.
				\end{proof}
				\subsubsection{Exactness}
				\begin{lemma}
					$Ker,Coker$ commutes with component operation $(-)_n:Ch(\mathcal{A})\rightarrow \mathcal{A}$.
				\end{lemma}
				\begin{proof}
					\begin{itemize}
						\item $Ker$:
						\item $Coker$:
					\end{itemize}
				\end{proof}
				\begin{proposition}
					The followings are equivalent:
					\begin{itemize}
						\item $0\rightarrow A\rightarrow B\rightarrow C\rightarrow0$ is exact in $Ch(\mathcal{A})$;
						\item $0\rightarrow A_n\rightarrow B_n\rightarrow C_n\rightarrow 0$ is exact for each $n\in \ZZ$ in $\mathcal{A}$;
					\end{itemize}
				\end{proposition}
				\begin{proof}
					By definition, omitted.	
				\end{proof}
				
				\begin{proposition}
					Suppose $A,B,C\in \mathcal{A},\mathcal{A}$ abelian, $0\rightarrow A\rightarrow B\rightarrow C\rightarrow 0$ is exact, then
					\begin{center}
						$0\rightarrow A\rightarrow B\rightarrow C\rightarrow 0$ splits$\Leftrightarrow B=A\oplus C$ 
					\end{center}
				\end{proposition}
				\begin{proof}
					Consider the following diagram
					\begin{center}
						\begin{tikzcd}
								0\arrow[r]&A\arrow[r,"f"]\arrow[rd]&B\arrow[d,shift right=2]\arrow[r,"g",shift right=2]&C\arrow[l,shift right=2,dashrightarrow]\arrow[ld]\arrow[r]&0\\
								&&A\oplus C\arrow[lu,shift left=2]\arrow[ru,shift right=2]\arrow[u,shift right=2,dashrightarrow]
						\end{tikzcd}
						Take $f\pi _A+h \pi_B:A\oplus C\rightarrow B$
					\end{center}
					$\pi_A,\pi_B$ satisfied by \ref{Homological-Algebra-Abelian-Category-Product-Equals-Coproduct}.
				\end{proof}
				
				\begin{proposition}
					(Left-Exactness of $Hom(-,X)$) Suppose $X\in \mathcal{A}$, then $Hom(-,X)$ is left-exact: For exact
					\begin{gather*}
						0\rightarrow A\rightarrow B\rightarrow C\rightarrow 0
					\end{gather*}
					We get the left-exact sequence
					\begin{gather*}
						0\rightarrow Hom(C,X)\rightarrow Hom(B,X)\rightarrow Hom(A,X)
					\end{gather*}
				\end{proposition}
				\begin{proof}
					The morphisms induced by the morphims above.
					\begin{itemize}
						\item $Hom(C,X)\rightarrow Hom(B,X)$ is injective: We have
						\begin{gather*}
							h\circ g=0,g\mbox{ is surjective}\Rightarrow h=0
						\end{gather*}
						\item $Im(f^*)=Ker(g^*)$: 
						\begin{itemize}
							\item $Im(g^*)\subseteq Ker(f^*)$, easy to verify, that for $h\in Hom(C,X)$,
							\begin{gather*}
								f^*g^*(h)=h\circ g\circ f=0
							\end{gather*}
							\item $Ker(f^*)\subseteq Im(g^*)$, we have for $h\in Hom(B,X)$,
							\begin{gather*}
								f^*(h)=h\circ f=0
							\end{gather*}
							Now, we can get the diagram
							\begin{center}
								\begin{tikzcd}
									0\arrow[r]&A\arrow[r,"f"]&B\arrow[r,"g"]\arrow[d,"h"]&C\arrow[r]\arrow[ld,dashrightarrow]&0\\
									&&X
								\end{tikzcd}
							\end{center}
							Notice that $Coker(f)=C$. Thus the morphism $C\rightarrow X$ is induced.
						\end{itemize}
					\end{itemize}
				\end{proof}
				Similarly
				\begin{proposition}
					(Left Exactness of $Hom(X,-)$) Suppose $X\in \mathcal{A}$, then $Hom(-,X)$ is right-exact:For exact
					\begin{gather*}
						0\rightarrow A\rightarrow B\rightarrow C\rightarrow 0
					\end{gather*}
					Wwe can get a right-exact sequence
					\begin{gather*}
						0\rightarrow Hom(X,A)\rightarrow Hom(X,B)\rightarrow Hom(X,C)
					\end{gather*}
				\end{proposition}
				\begin{proof}
					\begin{itemize}
						\item $f^*:Hom(X,A)\rightarrow Hom(X,B)$ is injective: We have
						\begin{gather*}
							f^*(h)=0\Rightarrow f\circ h=0\Rightarrow h=0
						\end{gather*}
						\item Similarly $Im(f^*)=Ker(g^*)$.
					\end{itemize}
				\end{proof}
				Furthermore, we have a longer exact sequence: Here
				\subsection{Homology Sequence}
					\begin{Def}
						Define the homology sequence, with objects:
						\begin{gather*}
							H_n(C):=Im(d_{n+1})/ker(d_n)
						\end{gather*}
						with the 
					\end{Def}
					\begin{proposition}
						$H_\bullet(-)$ is a morphism in $Ch(-)$.
					\end{proposition}
					\begin{proof}
						Omitted.
					\end{proof}
					\begin{Def}
						Define: A morphism $C\rightarrow D$ is a \textbf{quasi-isomorphism} if the induced map
						\begin{gather*}
							H_n(f):H(C)\rightarrow H(D)
						\end{gather*}
						is an isomorphism.
					\end{Def}
				\subsection{Operations on Chain Complexes}
				\subsubsection{Product, Direct Sum}
					\begin{Def}
						For a family of chain complexes $\{C_\alpha\}$ in $Ch(\mathcal{A})$, where $\mathcal{A}$ is abelian category, define:
						\begin{itemize}
							\item Product $(\prod A_\alpha,\prod d_\alpha)$, with
							\begin{gather*}
								\prod d_\alpha:\prod A_{\alpha,n}\rightarrow \prod A_{\alpha,n-1} 
							\end{gather*}
							\item Sum $(\oplus A_\alpha,\oplus d_\alpha)$, with	
							\begin{gather*}
								\oplus d_\alpha:\oplus A_{\alpha,n}\rightarrow \oplus A_{\alpha,n-1}
							\end{gather*}
						\end{itemize}
					\end{Def}
					\begin{proposition}
						The two operations commute with homology, i.e
						\begin{gather*}
							\oplus H_n(A_\alpha)\cong H_n(\oplus A_\alpha)\qquad \prod H_n(A_\alpha)=H_n(\prod A_\alpha)
						\end{gather*}
					\end{proposition}
					\begin{proof}
						内容...
					\end{proof}
				\subsubsection{Subcomplex,Quotient Complex}
					\begin{Def}
						A \textbf{subcomplex}
					\end{Def}
					Especially: In $Ab$ the quotient structure is defined as:
					\begin{Def}
						Given $A\hookrightarrow B$, then
						\begin{gather*}
							A/B=coker(A\rightarrow B)
						\end{gather*}
					\end{Def}
					Thus, define the quotient complex
					\begin{Def}
						Define: For subcomplex $(B_n)\hookrightarrow (A_n)$ define
						\begin{gather*}
							((A_n/B_n),(\bar{d}_n))
						\end{gather*}
						Induced by:
						\begin{center}
							\begin{tikzcd}
								&0\arrow[d]&0\arrow[d]&0\arrow[d]\\
								\dots\arrow[r]&B_{n-1}\arrow[r]\arrow[d]&\arrow[r]B_n\arrow[r]\arrow[d]&B_{n+1}\arrow[r]\arrow[d]&\dots\\
								\dots\arrow[r]&A_{n-1}\arrow[d]\arrow[r]&A_n\arrow[r]\arrow[d]&A_{n+1}\arrow[d]\arrow[r]&\dots\\
								\dots\arrow[r]&A_{n-1}/B_{n-1}\arrow[r]\arrow[d]&A_n/B_n\arrow[r]\arrow[d]&A_{n+1}/B_{n+1}\arrow[r]\arrow[d]&\dots\\
								&0&0&0
							\end{tikzcd}
						\end{center}
					\end{Def}
					The exactness is given by \underline{9-lemma}.
				\subsubsection{Tensor Product}
				\begin{Def}
					For $X,Y\in Ch_\bullet(A)$ write $X\otimes Y\in Ch_\bullet(A)$ for chain complex whose component is given by:
					\begin{gather*}
						(X\otimes Y)_n:=\oplus_{i+j=n} X_i\otimes_R Y_j
					\end{gather*}
					with $\partial^{x\otimes y}(x,y)=(\partial^Xx,y)+(-1)^{deg(X)}(x,\partial^Yy)$.
				\end{Def}
				
				\begin{example}
					Let $R$ be some ring and $I_\bullet\in C_\bullet(R-Mod)$ be the chain complex given by
					\begin{gather*}
						I_\bullet=[\dots\rightarrow 0\rightarrow 0\rightarrow R\stackrel{\partial_0^I}{\rightarrow} R\oplus R]
					\end{gather*}
					where $\partial_0^I=(-id,id)$. We may think of 
					\begin{gather*}
						I_0=R\oplus R\cong R[\{(0),(1)\}]\qquad I_1=R\cong R[(0\rightarrow 1)]
					\end{gather*}
					Thus,
					\begin{gather*}
						\partial_0^I:(0\rightarrow 1)\mapsto (1)-(0)
					\end{gather*}
					Write out the tensor product $S_\bullet:=I_\bullet\otimes I_\bullet$, by definition we have
					\begin{gather*}
						S_0=\\
						S_1=\\
						S_2=
					\end{gather*}
				\end{example} 
				
				\begin{proposition}
					For category $\mathcal{A}$, we have
					\begin{center}
						$Ch(\mathcal{A})$ is monoidal$\Leftrightarrow$ $\mathcal{A}$ is monoidal
					\end{center}
				\end{proposition}
				\begin{proof}
					内容...
				\end{proof}
				\subsubsection{Chain $[X,Y]$}
				\begin{Def}
					\begin{itemize}
						\item 
						Define the category $Ch_\bullet^{ub}(\mathcal{A}):=\{C\in Ch_\bullet(\mathcal{A})|C$ is unbounded$\}$
						\item Define:
						\begin{gather*}
							[X,Y]_n:=\prod_{i\in \ZZ} Hom_{RMod}(X_i,Y_{i+n})\\
							df:=d_Y\circ f-(-1)^nf\circ d_X \quad f\in [X,Y]_n
						\end{gather*}
					\end{itemize}
				\end{Def}
				Thus,
				\begin{gather*}
					[-,-]:Ch_\bullet(\mathcal{A})^{op}\times Ch_\bullet(\mathcal{A})\rightarrow Ch_\bullet^{ub}(\mathcal{A})
				\end{gather*}
				is the \textbf{inner hom} of the category of chain complexes.
				
				\begin{proposition}
					The collection of cycles of the internal $Hom[X,Y]_\bullet$ in degree $0$ coincides with the external hom functor
					\begin{gather*}
						Z_0([X,Y])=Hom_{Ch_\bullet^{ub}(X,Y)}
					\end{gather*}
				\end{proposition}
				\begin{proof}
					内容...
				\end{proof}
				
				\begin{proposition}
					The monoidal category $(Ch_\bullet,\otimes)$ is a closed monoidal category, the internal hom is the standard internal hom of chain complexes.
				\end{proposition}
				\begin{proof}
					内容...
				\end{proof}
			\subsection{Shifting and Translation}
				\begin{Def}
					Define the chain/cochain complex
					\begin{gather*}
						C[p]_n:=C_{n+p}\qquad C[p]^n=C^{n-p}
					\end{gather*}
				\end{Def}
				Then, define the differential map
				\begin{gather*}
					d:=
				\end{gather*}
				
			\subsection{Bicomplex}
				\begin{Def}
					A \textbf{double complex} or \textbf{bicomplex} in $\mathcal{A}$ is a family $\{C_{p,q}\}$ of object $\mathcal{A}$ together with maps
					\begin{gather*}
						d^h:C_{p,q}\rightarrow C_{p-1,q}\qquad d^v:C_{p,q}\rightarrow C_{p,q-1}
					\end{gather*}
					Such that: $d^h\circ d^h=d^v\circ d^v=d^vd^h+d^hd^v=0$. Lattice
					\begin{center}
						{\small	\begin{tikzcd}
								&\dots\arrow[d]&\dots\arrow[d]&\dots\arrow[d]\\
								\dots\arrow[r]&C_{p-1,q+1}\arrow[r]\arrow[d]&C_{p-1,q}\arrow[r]\arrow[d]&C_{p-1,q-1}\arrow[r]\arrow[d]&\dots\\
								\dots\arrow[r]&C_{p,q+1}\arrow[r]\arrow[d]&C_{p,q}\arrow[d]\arrow[r]\arrow[d]&C_{p,q-1}\arrow[r]&\dots\\
								\dots\arrow[r]&C_{p+1,q+1}\arrow[r]\arrow[d]&C_{p+1,q}\arrow[r]\arrow[d]&C_{p+1,q-1}\arrow[d]\arrow[r]&\dots\\
								&\dots&\dots&\dots
						\end{tikzcd}}
					\end{center}
				\end{Def}
			\subsubsection{Tot Complexes}
			\begin{Def}
				Define:
				\begin{gather*}
					Tot^{\prod}_n:= \prod_{p+q=n} C_{p,q}\qquad  Tot^{\oplus}_n=\oplus_{p+q=n}C_{p,q}
				\end{gather*}
				With the differential
				\begin{gather*}
					d:Tot^{\prod}_n\rightarrow Tot^{\prod}_{n-1}\\
					d:Tot^{\oplus}_n\rightarrow Tot^{\oplus}_{n-1}
				\end{gather*}
				such that $d\circ d=0$.
			\end{Def}
			
			\begin{proposition}
				$C$ is a bounded double complex with exact rows/columns$\Rightarrow Tot(C)$ is acyclic.
			\end{proposition}
			\begin{proof}
				内容...
			\end{proof}
			\subsection{Truncations,Translations}
			\subsubsection{Truncations}
			\begin{Def} If $C$ is a chain complex, and $n$ is an integer. We let $\tau_{\geq n} \mathcal{C}$ denote the subcomplex of $\mathcal{C}$ defined by
				\begin{gather*}
					(\tau_{\geq n} C)_i=\begin{cases}
						0&i<n\\
						Z_n&i=n\\
						C_i&i>n
					\end{cases}
				\end{gather*}
				Conversely, the quotient complex $\tau_{<n} C=C/(\tau_{\geq n} C)$
			\end{Def}
			\begin{proposition}
				Chearly that:
				\begin{gather*}
					H_i(\tau_{\geq n}C)=\begin{cases}
						0&i<n\\
						H_i(C)&i\geq n
					\end{cases}
				\end{gather*}
			\end{proposition}
			\begin{proof}
				内容...
			\end{proof}
			
			
			\subsubsection{Translations}
			\begin{Def}
				Let $C$ be a complex. Then, define
				\begin{gather*}
					C[p]_n=C_{n+p}\qquad C[p]^n=C^{n-p}
				\end{gather*}
			\end{Def}
			More information will be given here:
		\section{Long Exact Sequence}
			\subsection{9-Lemma}
			\begin{proposition}
				(9-Lemma) Suppose the following commutative diagram
				\begin{center}
					\begin{tikzcd}
						&0&0&0\\
						0&A'&B'&C'&0\\
						0&A&B&C&0\\
						0&A''&B''&C''&0\\
						&0&0&0
					\end{tikzcd}
				\end{center}
				In abelian category $\mathcal{A}$, then:
				\begin{itemize}
					\item If the bottom two rows are exact, so is the top;
					\item If the top two rows are exact, so is the bottom row;
					\item If the top and bottom rows are exact, and the composite $A\rightarrow C$ is zero, the middle row is also exact;
				\end{itemize}
			\end{proposition}
			\begin{proof}
				
			\end{proof}
			\subsubsection{Snake's Lemma}
			\begin{corollary}
				(Snake's Lemma) Consider the following commutative diagram
				\begin{center}
					\begin{tikzcd}
						&A\arrow[r]\arrow[d]&B\arrow[r]\arrow[d]&C\arrow[r]\arrow[d]&0\\
						0\arrow[r]&A'\arrow[r]&B'\arrow[r]&C'
					\end{tikzcd}
				\end{center}
				There exists an exact sequence:
				\begin{center}
					\begin{tikzcd}
						0\arrow[r]&ker(f)\arrow[r]&ker(g)\arrow[r]&ker(h)\arrow[r]&Coker(f)\arrow[r]&Coker(g)\arrow[r]&Coker(h)\arrow[r]&0
					\end{tikzcd}
				\end{center}
			\end{corollary}
			\begin{proof}
				内容...
			\end{proof}
			\subsubsection{5-Lemma}
			\begin{corollary}
				(5-Lemma) Consider the following diagram
				\begin{center}
					\begin{tikzcd}
						0\arrow[r]&A_1\arrow[r]\arrow[d]&A_2\arrow[r]\arrow[d]&A_3\arrow[r]\arrow[d]&A_4\arrow[r]\arrow[d]&A_5\arrow[r]\arrow[d]&0\\
						0\arrow[r]&B_1\arrow[r]&B_2\arrow[r]&B_3\arrow[r]&B_4\arrow[r]&B_5\arrow[r]&0
					\end{tikzcd}
				\end{center}
				with following properties
				\begin{itemize}
					\item If$\Rightarrow$
					\item If$\Rightarrow$
					\item If $\Rightarrow$
				\end{itemize}
			\end{corollary}
			\begin{proof}
				内容...
			\end{proof}
			\subsubsection{3-Lemma}
			
			\begin{corollary}
				(3-Lemma)For the diagram
				\begin{center}
					\begin{tikzcd}
						0&A_1&B_1&C_1&0\\
						0&A_2&B_2&C_2&0
					\end{tikzcd}
				\end{center}
			\end{corollary}
			\begin{proof}
				内容...
			\end{proof}
			\subsubsection{Long Exact Sequence}
			\begin{theorem}
				\begin{itemize}
					\item Consider the following short exact sequence in $\mathcal{A}$:
					\begin{gather*}
						0\rightarrow A\rightarrow B\rightarrow C\rightarrow 0
					\end{gather*}
					be an exact sequence, then we have a long exact sequence 
					\begin{gather*}
						\dots\rightarrow H_n(A)\rightarrow H_n(B)\rightarrow H_n(C)\stackrel{\delta}{\rightarrow} H_{n+1}(A)\rightarrow\dots
					\end{gather*}
					\item Conversely, $H^{n+1}(A)\rightarrow H^n(C)$.
				\end{itemize}
				And in this sense, $H_\bullet$ is a covariant functor.
			\end{theorem}
			\begin{proof}
				\begin{itemize}
					\item 
					\item 
				\end{itemize}
			\end{proof}
			\subsection{Exact Triangle}
				\begin{Def}
					内容...
				\end{Def}
				
				\begin{corollary}
					The category of chain complexes $Ch(\mathcal{A})$ is a \textbf{triangulated category}.
				\end{corollary}
				\begin{proof}
					内容...
				\end{proof}
		\section{Chain Homotopies}
			\begin{Def}
				A category 
				\begin{itemize}
					\item has \textbf{enough projective} if for every obeject $X$ if there is a projective object $Q$ equipped with an epimorphism $Q\rightarrow X$;
					\item has \textbf{enough injectives} iff for every object $X$ there is an injective $Q$ equipped with monomorphism $X\rightarrow P$. 
				\end{itemize}
			\end{Def}
			\begin{proposition}
				Assuming the axiom of choice, the category $R-Mod$ has enough projective module.
			\end{proposition}
			\begin{proof}
				内容...
			\end{proof}
			\subsection{Introduction}
				Consider the following vector space decomposition over field:
				\begin{lemma}
					Every exact sequence in $Vect_k$ splits.
				\end{lemma}
				\begin{gather*}
					C_n=Z_n\oplus B_n';Z_n=Im(\partial_{n+1}),B_n':=C_n/Z_n=Coker(Z_n\rightarrow C_n)\cong B_{n-1}\\
					Z_n=B_n\oplus H_n',H_n':=Z_n/B_n=H_n(C)
				\end{gather*}
				\begin{center}
					{\tiny \begin{tikzcd}
							&\dots\arrow[d]\\
							&C_{n+1}\arrow[d,"\partial_{n+1}"]&0\arrow[d]\\
							0\arrow[r]&Im(\partial_{n+1})=B_{n}\arrow[r]\arrow[u,dashrightarrow,shift left= 2]&Ker(\partial_n)=Z_{n}\arrow[r]\arrow[d]\arrow[l,shift left=2,dashrightarrow]&H_n(C_\bullet)\arrow[r]&0\\
							&&C_{n}\arrow[d,"\partial_{n}"]\arrow[u,shift left=2,dashrightarrow]\\
							&0\arrow[r]&Im(\partial_n)\arrow[r]=B_{n-1}\arrow[d]&Ker(\partial_{n-1})=Z_{n-1}\arrow[r]\arrow[d]&H_{n-1}(C_\bullet)\arrow[r]&0\\
							&&0&C_{n-1}\arrow[d]\\
							&&&\dots
					\end{tikzcd}}
				\end{center}
				Therefore, we can form thecompositions
				\begin{gather*}
					C_n\stackrel{\pi}{\rightarrow}Z_n\stackrel{\pi}{\rightarrow} B_{n}\cong B_{n+1}'\subseteq C_{n+1}
				\end{gather*}
				By the subspace $C_{n+1}/Z_{n+1}\hookrightarrow C_{n+1}$. Define the \textbf{split map} $s_n:C_n\rightarrow C_{n+1}$. We have
				\begin{gather*}
					\partial_{n+1}s_n\partial_{n+1}=\partial_{n+1}
				\end{gather*}
				Verify:
				\begin{gather*}
					内容...
				\end{gather*}
				And:
				\begin{itemize}
					\item $\partial s$:
					\item $s\partial$:
				\end{itemize}
				\subsubsection{Split}
				\begin{Def}
					A complex $C$ is called split if there exists some $s_n:C_{n}\rightarrow C_{n+1}$ such that
					\begin{gather*}
						d=dsd
					\end{gather*}
					such $s$ is called \textbf{splitting map}.
					
					Especially: If $C$ is exact, then we say that $C$ is \textbf{split exact}.
				\end{Def}
				\subsubsection{Chain Contruction}
					\begin{Def}
						\begin{itemize}
							\item We say that a chain map $f:C\rightarrow D$ is \textbf{null homotopic} if there are maps $s_n:C_n\rightarrow D_{n+1}$ such that
							\begin{gather*}
								f=ds+sd
							\end{gather*}
							\item The maps $\{s_n\}$ are called a \textbf{chain constraction} of $f$.
						\end{itemize}
					\end{Def}
					\begin{proposition}
						For chain complex $C_*$, $C_*$ splits exact$\Leftrightarrow id_{C_*}$ is null homotopic.
					\end{proposition}
					\begin{proof}
						\begin{itemize}
							\item $\Rightarrow$ Then $id_{C_*}=ds+sd$.  Now, we have
							\begin{gather*}
								d=d(ds+sd)=dsd
							\end{gather*}
							\item $\Leftarrow$ If $id_{C_*}$ is null homotopic then there exists some $s_n$ such that
							\begin{gather*}
								id_{C_n}=d_{n+1}s_n+s_{n-1}d_n
							\end{gather*}
							
							(Exactness) Suppose that $a\in C_n$, $d_n(a)=0,a\in Ker(d_n)$. Then,
							\begin{gather*}
								\Rightarrow s_{n-1}d_n(a)=0
								\\ \Rightarrow a=d_{n+1}s_n(a)\\
								\Rightarrow a\in Im(d_{n+1})\subseteq ker(d_n) \\
								d^2=0\Rightarrow Im(d_{n+1})=Ker(d_n)
							\end{gather*}
							
							(Splitting) 
							\begin{gather*}
								d=d(ds+sd)=dsd
							\end{gather*}
						\end{itemize}
					\end{proof}
				\subsubsection{Chain Homotopies}
					\begin{Def}
						\begin{itemize}
							\item $f:C\rightarrow D$ is \textbf{null homotopic} if
							\begin{gather*}
								f=ds+sd
							\end{gather*}
							\item Consider two chain maps 
							\begin{gather*}
								f,g:C\rightarrow D
							\end{gather*}
							Then, $f$ and $g$ are \textbf{homotopic} iff $f-g$ is null homotopic. i.e.
							\begin{gather*}
								f-g=sd+ds
							\end{gather*}
							The maps $(s_n)$ are called \textbf{chain homotopy from} $C$ to $D$. 
							\item $f:C\rightarrow D$ is a \textbf{homotopy equivalence}
						\end{itemize}
					\end{Def}
					
					\begin{lemma}
						(Chain Homotopy Invariance for $H_\bullet$)  
						\begin{itemize}
							\item 
							\item 
						\end{itemize} 
					\end{lemma}
					\begin{proof}
						内容...
					\end{proof}
					
					\begin{proposition}
						\begin{itemize}
							\item 	If $f:C\rightarrow D$ is null homotopic, then every map $f_*:H_n(C)\rightarrow H_n(D)$ is zero;
							\item If $f\simeq g$, then $H_n(f)=H_n(g):H_n(C)\rightarrow H_n(D)$
						\end{itemize}
					\end{proposition}
					\begin{proof}
						内容...
					\end{proof}
				\subsubsection{Quotient Category}
					\begin{proposition}
						Let $K$ be a quotient category of $Ch$.Then, the functor
						\begin{gather*}
							H_n:Ch(\mathcal{A})\rightarrow \mathcal{A}
						\end{gather*}
						will factors through:
						\begin{itemize}
							\item 
							\item 
							\item 
							\item 
						\end{itemize}
					\end{proposition}
					\begin{proof}
						内容...
					\end{proof}
			\subsection{Resolution}
				
		\section{Homology Exact Sequence}
			\begin{Def}
				An \textbf{exact sqeunce} in $\mathcal{A}$ is a chain complex $C_\bullet$ in $\mathcal{A}$ with vanishing chain homology in each degree
				\begin{gather*}
					\forall n\in \NN,H_n(C)=0
				\end{gather*}
				
				Especially, a \textbf{short exact sequence} is an exact sequence of the form
				\begin{gather*}
					\dots\rightarrow0\rightarrow A\rightarrow B\rightarrow C\rightarrow 0\rightarrow \dots
				\end{gather*}
				Mostly written as
				\begin{gather*}
					0\rightarrow A\rightarrow B\rightarrow C\rightarrow 0
				\end{gather*}
			\end{Def}
			\subsection{Some of Skills in Chain }
				\subsubsection{$3\times 3$-Lemma}
					\begin{proposition}
						内容...
					\end{proposition}
					\begin{proof}
						内容...
					\end{proof}
				\subsubsection{Snake Lemma}
				\subsubsection{5-Lemma}
				\subsubsection{Addenbum}
			\subsection{Existence of $\partial$ in Exact Chain Complexes}
			So, the following diagram commutes:
			\begin{center}
				\begin{tikzcd}
					H_*(A)&&H_*(B)\\
					&H_*(C)
				\end{tikzcd}\qquad(The \textbf{exact triangle})
			\end{center}
			
			\begin{proposition}
				
			\end{proposition}
			\begin{proof}
				内容...
			\end{proof}
			\subsection{Resolution}
			\begin{proposition}
				Suppose $A_*$ be a resoltion of $M$, then
				\begin{gather*}
					H_n(A_*)=\begin{cases}
						M&n=0\\
						0&n\in \ZZ-\{0\}
					\end{cases}
				\end{gather*}
			\end{proposition}
			\begin{proof}
				Omitted.
			\end{proof}
			\subsection{Functor $Ext^1$}
			\begin{Def}
				Select an injective resolution of $X$:
				\begin{gather*}
					0\rightarrow X\rightarrow I^0\rightarrow I^1\rightarrow I^2\rightarrow\dots
				\end{gather*}
				Apply $H(A,-)$ and deleting $Hom(A,X)$ we get
				\begin{gather*}
					0\rightarrow Hom(A,I^0)\stackrel{d^0}{\rightarrow} Hom(A,I^1)\stackrel{d^1}{\rightarrow} Hom(A,I^2)\rightarrow \dots
				\end{gather*}
				Define:
				\begin{gather*}
					Ext^1(A,X):=ker(d^1)/im(d^0)
				\end{gather*}
			\end{Def}
			We have the following isomophic construction via projective resolution:
			\begin{proposition}
				In in the category of $\ZZ$-module/abelian groups,we have the exact sequence as followings:
				\begin{center}
					\begin{tikzcd}
						0\arrow[r]&Hom(C,X)\arrow[r]&Hom(B,X)\arrow[r]&Hom(A,X)\arrow[dll, {name=D, offset=-1}, out=0, in=180]\\
						&Ext^1(C,X)\arrow[r]&Ext^1(B,X)\arrow[r]&Ext^1(A,X)\arrow[r]&0
					\end{tikzcd}
				\end{center}
			\end{proposition}
			\begin{proof}
				内容...
			\end{proof}
			In General: In $\mathcal Ab$-category we have the long exact sequence (here ) 
			\subsubsection{Extension and Ext}
				Ref: Weibel 3.4
				\begin{Def}
					\begin{itemize}
						\item An \textbf{extension} of $A$ by $B$ is an exact sequence
						\begin{gather*}
							0\rightarrow A\rightarrow X\rightarrow B\rightarrow 0
						\end{gather*}
						\item Two extensions of $A$ by $B$ are \textbf{equivalent} if
						\begin{center}
							\begin{tikzcd}
								0\arrow[r]&B\arrow[r]\arrow[d,equal]&X\arrow[r]\arrow[d]&A\arrow[r]\arrow[d,equal]&0\\
								0\arrow[r]&B\arrow[r]&X\arrow[r]&A\arrow[r]&0
							\end{tikzcd}
						\end{center}
					\end{itemize}
				\end{Def}
				\begin{example}
					(Extension of $\ZZ/p$) Consider the following equivalent class of split extensions:
					\begin{center}
						\begin{tikzcd}
							0\arrow[r]&\ZZ/p\arrow[r,"p"]&\ZZ/p^2\arrow[r,"i"]&\ZZ/p\arrow[r]&0
						\end{tikzcd},i=1,\dots,p-1
					\end{center}
				\end{example}
				\begin{lemma}
					$Ext^1(A,B)=0\Leftrightarrow $ Every extension of $A$ by $B$ splits.
				\end{lemma}
				\begin{proof}
					内容...
				\end{proof}
				\begin{theorem}
					内容...
				\end{theorem}
				\begin{proof}
					内容...
				\end{proof}
				\begin{Def}
					(Baer Sum)
				\end{Def}
		\section{Mapping Cones and Cylinders}
			\subsection{Mapping Cones}
			\begin{Def}
				Let $f:B\rightarrow C$ be a map of chain complexes. The \textbf{mapping cone} of $f$ is the chain complex $cone(f)$ given by:
				\begin{gather*}
					[cone(f)]_n:=B_{n-1}\oplus C_n
				\end{gather*}
				With differentials
				\begin{center}
					$d(b,c)=(-d(b),d(c)-f(b))\quad\begin{pmatrix}
						-d_B&0\\
						-f&+d_C
					\end{pmatrix}:$\qquad{\tiny \begin{tikzcd}
					B_{n-1}\arrow[rr,"-"]\arrow[rrdd,"-"]&&B_{n-2}\\
					\oplus&&\oplus\\
					C_n\arrow[rr,"+"]&&C_{n-1}
				\end{tikzcd}}
				\end{center}
				Then, we form a exact sequence of double complexes:
				\begin{center}
					\begin{tikzcd}
						0\arrow[r,"i"]&C\arrow[r]&D=cone(f)\arrow[r,"p"]&B[-1]\arrow[r]&0
					\end{tikzcd}
				\end{center}
			\end{Def}
			
			Especially $cone(C):=cone(id_C)$
			
			\textbf{Remark} This is a double complex: Verify
			\begin{center}
				\begin{tikzcd}
					C_{n-1}\arrow[r]\arrow[d]&B_{n-2}\oplus C_{n-1}\arrow[r]\arrow[d]&B_{n-2}\arrow[d]\\
					C_n\arrow[r]&B_{n-1}\oplus C_n\arrow[r]&B_{n-1}
				\end{tikzcd}
			\end{center}
			We have $d(i(c))=i(d(c))=(0,d(c));c\in C_{n-1};p(d(b,c))=d(p(b,c))=d(c)$
			
			And exactness by definition.
			\begin{proposition}
				$f:C\rightarrow D$ is \textbf{null homomtopic} iff if $f$ extends $(-s,f):cone(C)\rightarrow D$.
			\end{proposition}
			\begin{proof}
				内容...
			\end{proof}
			
			\begin{proposition}
				With the exact sequence above, we have a long exact sequence
				\begin{gather*}
					\dots\rightarrow H_{n+1}(cone(f))\rightarrow H_n(B)\rightarrow H_n(C)\rightarrow H_n(cone(f))\rightarrow \dots
				\end{gather*}
				with $\partial=f_*:H_*(B[-1])\rightarrow H_*(C)$
			\end{proposition}
			\begin{proof}
				内容...
			\end{proof}
			
			\begin{corollary}
				$f:B\rightarrow C$ is a quasi-isomorphism$\Leftrightarrow$ The mapping cone of $cone(f)$ is exact.
			\end{corollary}
			\begin{proof}
				内容...
			\end{proof}
			\subsection{Mapping Cylinder}
			\begin{Def}
				Define: For $f:B\rightarrow C$, the \textbf{mapping cylinder}
				\begin{gather*}
					\begin{cases}
						cyl(f):=B_n\oplus B_{n-1}\oplus C_n\\
						d(b,b',c):=(d(b)+b',-d(b'),d(c)-f(b'))
					\end{cases}\\ \begin{pmatrix}
						d_B&id_B&0\\
						0&-d_B&0\\
						0&-f&d_C
					\end{pmatrix}\Rightarrow \begin{pmatrix}
					d_B&id_B&0\\
					0&-d_B&0\\
					0&-f&d_C
					\end{pmatrix}^2=\begin{pmatrix}
						d_B^2&d_B-d_B&0\\
						0&d_B^2&0\\
						0&fd_B-d_Cf&d_C^2
					\end{pmatrix}=0
				\end{gather*}
				The diagram as following:
				\begin{center}
					{\tiny \begin{tikzcd}
							B_n\arrow[rr,"+"]&&B_{n-1}\\
							\oplus&&\oplus\\
							B_{n-1}\arrow[rr,"-"]\arrow[rruu,"+"]\arrow[rrdd,"-"]&&B_{n-2}\\
							\oplus&&\oplus\\
							C_n\arrow[rr,"+"]&&C_{n-1}
					\end{tikzcd}}
				\end{center}
			\end{Def}
			
			\begin{lemma}
				Let $cyl(C):=cyl(id_C)$. Then, it has 
				\begin{gather*}
					C_n\oplus C_{n-1}\oplus C_n
				\end{gather*}
				in degree $n$. Then, the two chain maps $f,g:C\rightarrow D$ are chain homotopic$\Leftrightarrow$ They extends to a map
				\begin{gather*}
					(f,s,g):cyl(C)\rightarrow D
				\end{gather*}
			\end{lemma}
			\begin{proof}
				\begin{itemize}
					\item $\Rightarrow$ If $f-g=sd+ds$ then we have
					\begin{gather*}
						(f,s,g):cyl(C)\rightarrow D\qquad (a,b,c)\mapsto f_n(a)+s_{n-1}(b)+g_n(c)
					\end{gather*}
					Then, we have
					\begin{gather*}
						(f,s,g)\circ d_{Cyl(C)}=\begin{pmatrix}
							f&s&g
						\end{pmatrix}\begin{pmatrix}
							d&1&0\\
							0&-d&0\\
							0&-1&d
						\end{pmatrix}\begin{pmatrix}
							1\\
							1\\
							1
						\end{pmatrix}=fd+f-ds-g+gd=(s+f+g)d=d(s+f+g)
					\end{gather*}
					\item $\Leftarrow$ Suppose that $(f,s,g)$ is a chain map, then
					\begin{gather*}
						(f,s,g)\circ d_{Cyl(C)}=fd+f-ds-g+gd=fd+gd+sd\Rightarrow f-g=ds+ds
					\end{gather*}
				\end{itemize}
			\end{proof}
			
			\begin{lemma}
				The subcomplex of elements $(0,0,c)$ is isomorphic to $C$. 
				
				The corresponding inclusion 
				\begin{gather*}
					\alpha:C\rightarrow cyl(f)
				\end{gather*}
				
				is a \underline{quasi-isomorphism}.
			\end{lemma}
			\begin{proof}
				Firstly, this is a subcomplex. We have
				\begin{gather*}
					d(0,0,c)=(0,0,d(c))
				\end{gather*}
				
				Then, $\alpha:C\rightarrow Cyl(f),c\mapsto (0,0,c)$, then
				\begin{gather*}
					H_n(\alpha):H_n(C)\rightarrow H_{n}(cyl(f))
				\end{gather*}
				
				is a quasi-isomorphism: The quotient $cyl(f)/\alpha(C)=cone(-id_B)$. 
				\begin{gather*}
					内容...
				\end{gather*}
				So, it is null-homotopic, thus follows the long exact homology sequence
				\begin{gather*}
					0\rightarrow C\stackrel{\alpha}{\rightarrow} cyl(f)\rightarrow cone(-id_B)\rightarrow 0
				\end{gather*}
			\end{proof}
			
			\begin{proposition}
				As above. $\beta:(b,b',c)\mapsto f(b)+c,cyl(f)\rightarrow C$ defines a chain map such that
				\begin{gather*}
					\beta\alpha=id_C
				\end{gather*}
				
				And, $s(b,b',c)\rightarrow (0,b,0)$ defines a chain homotopy from the identity of $cyl(f)$ to $\alpha\beta$.
				
				$\alpha$ is in fact a chain homotopy equivalence between $C$ and $cyl(f)$.
			\end{proposition}
			\begin{proof}
				Consider the composition
				\begin{gather*}
					内容...
				\end{gather*}
			\end{proof}
			\subsubsection{Exact Sequence between }
		\section{Euler Characteristics}
			\subsection{Betti Numbers}
			\subsection{Euler Characteristics}
			\subsubsection{On Complex of Vector Spaces}
			\begin{Def}
				Let $R$ be a commutative ring, $V$ be a chain complex of $R$-modules:
				
				If each $V_n$ is \textbf{finitely generated and projective}. Then, the \textbf{Euler characteristic} of $V$ is the alternating sum of their ranks, if this is finite
				\begin{gather*}
					\chi(V)=\sum (-1)^n rank_R(V_n)
				\end{gather*}
		
			\end{Def}
				\subsubsection{On the Chain Complexes}
			\begin{theorem}
				(Euler-Poincare) 
				\begin{gather*}
					\chi(C_\bullet)=\sum (-1)^i dim(H_i(C))
				\end{gather*}
			\end{theorem}
			\begin{proof}
				内容...
			\end{proof}
			\subsection{Euler Characteristic of a Topological Space}
			\subsection{From Euler Characteristics to $K_0$ Groups}
	\chapter{Derived Functors}
		\section{$\delta$-Functors}
			\begin{Def}
				A (covariant) homological/cohomological $\delta$-functor between $\mathcal{A}$ and $\mathcal{B}$ is a collection of additive functors $T_n:\mathcal{A}\rightarrow \mathcal{B}$ and
				\begin{gather*}
					\delta_n:T_n(C)\rightarrow T_{n-1}(A)\\
					\delta^n:T^n(C)\rightarrow T^{n+1}(A)
				\end{gather*}
				defined for each short exact sequence $0\rightarrow A\rightarrow B\rightarrow C\rightarrow 0$ in $\mathcal{A}$. Here we make the convention that $T^n=T_n=0,n<0$. 
			\end{Def}
			The two conditions are imposed:
			\begin{itemize}
				\item Long exact sequence
				\begin{gather*}
					\dots\rightarrow T_{n+1}(C)\stackrel{\delta}{\rightarrow} T_n(A)\rightarrow T_n(B)\rightarrow T_n(C)\stackrel{\delta}{\rightarrow}T_{n-1}(A)\rightarrow \dots\\
					\dots\rightarrow T^{n-1}(C)\stackrel{\delta}{\rightarrow} T^n(A)\rightarrow T^n(B)\rightarrow T^n(C)\stackrel{\delta}{\rightarrow}T^{n+1}(A)\rightarrow \dots
				\end{gather*}
				\item For two exact $0\rightarrow A\rightarrow B\rightarrow C\rightarrow 0$ and $0\rightarrow A'\rightarrow B'\rightarrow C'\rightarrow 0$, we have commutes squares
				\begin{center}
					\begin{tikzcd}
						T_n(C')\arrow[d]\arrow[r,"\delta"]&T_{n-1}(A')\arrow[d]\\
						T_n(C)\arrow[r,"\delta"]&T_{n-1}(A)
					\end{tikzcd}
				\end{center}
			\end{itemize}
			\begin{example}
				Consider the functors 
				\begin{gather*}
					H_*:Ch_{\geq 0}(\mathcal{A})\rightarrow\mathcal{A}\\
					H^*:Ch^{\geq 0}(\mathcal{A})\rightarrow \mathcal{A} 
				\end{gather*}
				$H_*$ gives a homological $\delta$-functor, and $H^*$ gives a cohomological $\delta$-functor. See here:
			\end{example}
			\begin{proposition}
				Let $\mathcal{S}$ be a category of short exact sequences
				\begin{gather*}
					0\rightarrow A\rightarrow B\rightarrow C\rightarrow 0\qquad(*)
				\end{gather*}
				in $\mathcal{A}$. Then, $\delta_i$ is a natural transformation form the functor sending from the functor sending $(*)$ to $T_i(C)$, to the functor sending $(*)$ to $T_{i-1}(A)$.
			\end{proposition}
			\begin{proof}
				Consider the two functors on on the category of short exact sequences $\mathcal{S}$ to $\mathcal{B}$:
				\begin{gather*}
					\mathcal{F}:(0\rightarrow A\rightarrow B\rightarrow C\rightarrow 0)\mapsto T_i(C)\\
					\mathcal{G}:(0\rightarrow A\rightarrow B\rightarrow C\rightarrow 0)\mapsto T_{i-1}(A)
				\end{gather*} 
				
				Thus, the natural transformation
				\begin{center}
					\begin{tikzcd}
						\mathcal{S}\arrow[rr,bend left=20,"\mathcal{F}"]\arrow[rr,bend right=20,"\mathcal{G}",swap]&{\tiny \Downarrow}&\mathcal{B}
					\end{tikzcd}
				\end{center}
				
				We could consider the following example: 
				\begin{center}
					\begin{tikzcd}
						(0\rightarrow A\rightarrow B\rightarrow C\rightarrow 0)\arrow[d,"{(f,g,h)}"]&T_i(C)\arrow[r]\arrow[d,"T_i(h)"]&T_{i-1}(A)\arrow[d,"T_{i-1}(f)"]\\
						(0\rightarrow A'\rightarrow B'\rightarrow C'\rightarrow 0)&T_i(C')\arrow[r]&T_{i-1}(A')
					\end{tikzcd}
				\end{center}
			\end{proof}
			
			\begin{example}
				(p-torsion) If $p$ is an integer,the functors $T_0(A):=A/pA$ and
				\begin{gather*}
					T_1(A)={}_pA=\{a\in A|pa=0\}
				\end{gather*}
				fit together to form a homological $\delta$-functor, or a cohomological $\delta$-functor from $Ab$ to $Ab$: Apply the snake lemma to
				\begin{center}
					\begin{tikzcd}
						0\arrow[r]&A\arrow[r]\arrow[d,"p"]&B\arrow[r]\arrow[d,"p"]&C\arrow[r]\arrow[d]&0\\
						0\arrow[r]&A\arrow[r]&B\arrow[r]&C\arrow[r]&0
					\end{tikzcd}
				\end{center}
				We get the exact sequence:
				\begin{gather*}
					0\rightarrow{ }_p A\rightarrow { }_p B\rightarrow { }_p C\stackrel{\delta}{\rightarrow} A/pA\rightarrow B/pB\rightarrow C/pC\rightarrow 0
				\end{gather*}
			\end{example}
			\begin{example}
				(Generalization) Similarly, for a ring $R$ with element $r\in R$, then
				\begin{gather*}
					T_0(M):=M/rM\qquad T_1:={ }_r M
				\end{gather*}
				
				And then, we will build a homological/cohomological functor $\delta:R-Mod\rightarrow {\mathcal Ab}$,
				\begin{gather*}
					T_n(M)=Tor_n^R(R/r,M)
				\end{gather*}
				
				Especially, if $r$ is nonzerodivisor, equivalently to say is 
				\begin{gather*}
					{ }_r R:=\{s\in R|rs=0\}=0
				\end{gather*}
				
				Then, we have
				\begin{gather*}
					T_1^R(M)=Tor_1^R(R/r,M)={ }_r M;Tor_n(R/r,M)=0,n\geq 2
				\end{gather*}
				
				And:
				\begin{gather*}
					Tor_1^R(R/r,M)={ }_r M/\{sm|rs=0,r\in R,m\in M\}
				\end{gather*}
				
				In general, ${ }_r R\not=0$.
			\end{example}
			
			\subsection{Morphism of $\delta$-Functors}
			\begin{Def}
				Suppose $S,T$ be two collection of $\delta$-funtors. A \textbf{morphism of $\delta$-functors} is a family of natural transformations between these functors, which commute with $\delta$. i.e. For $0\rightarrow A\rightarrow B\rightarrow C\rightarrow 0$ we have
				\begin{center}
					\begin{tikzcd}
						S_n(C)\arrow[r,"\delta_n"]\arrow[d,"f_n(C)"]&S_{n-1}(A)\arrow[d,"f_{n-1}(A)"]\\
						T_n(C)\arrow[r,"\delta_n"]&T_{n-1}(A)
					\end{tikzcd}
				\end{center}
			\end{Def}
			And
			\begin{Def}
				\begin{itemize}
					\item A homological functor $T$ is \textbf{universal}, if: Given
					\begin{itemize}
						\item Another $\delta$-Functor $S$;
						\item A morphism $f_0:S_0\rightarrow T_0$
					\end{itemize}
					There exists a unique morphism $\{f_n:S_n\rightarrow T_n\}$ of $\delta$ functor extends $f_0$;
					\item Dual Definition \textbf{couniversal} for extension $f_0:T_0\rightarrow S_0$ to $\{f_n:T_n\rightarrow S_n\}$.
				\end{itemize}
			\end{Def}
			
			\begin{proposition}
				Functor 
				\begin{itemize}
					\item $H_*(-):Ch_{\geq 0}(\mathcal{A})\rightarrow \mathcal{A}$ is universal $\delta$-functor;
					\item $H^*(-):Ch^{\geq 0}(\mathcal{A})\rightarrow \mathcal{A}$ is universal $\delta$-functor;
				\end{itemize}
			\end{proposition}
			\begin{proof}
				内容...
			\end{proof}
			\begin{proposition}
				Consider the following functor: Suppose $F:\mathcal{A}\rightarrow \mathcal{B}$ is exact, then the following functor:
				\begin{gather*}
					T_0=F\qquad T_n=0,n\not=0
				\end{gather*}
				Then, $T=\{T_n\}$ is a universal $\delta$-functor.
			\end{proposition}
			\begin{proof}
				内容...
			\end{proof}
		
		\section{Left Derived Functors}
			\begin{Def}
				Let $F:\mathcal{A}\rightarrow \mathcal{B}$ be a \underline{right exact functor} between two abelian categories. 
				
				
				If $\mathcal{A}$ has enough projectives, we can construct the \textbf{left derived functors} $L_iF,i\geq 0$ of $F$ as follows: If $A\in \mathcal{A}$, choose a projective resolution $P\rightarrow A$ and define
				\begin{gather*}
					L_iF(A):=H_i(F(P))
				\end{gather*}
			\end{Def}
			
			Notice
			
			\begin{center}
				\begin{tikzcd}
					\dots\arrow[r]& P_2\arrow[r]\arrow[d]&P_1\arrow[r]\arrow[d]&P_0\arrow[d]\\
					\dots\arrow[r]&0\arrow[r]&0\arrow[r]&A
				\end{tikzcd}$\Rightarrow$ Exact $ P_1\rightarrow P_0\rightarrow A\Rightarrow$ Exact $ F(P_1)\rightarrow F(P_0)\rightarrow F(A)\rightarrow 0$
			\end{center}
			
			Consider the exact sequence $F(P_1)\rightarrow F(P_0)\rightarrow F(A)\rightarrow 0$, we have
			\begin{gather*}
				L_0F(A)=H_0(F(P))=\frac{ker(F(A)\rightarrow 0)}{Im(F(P_0)\rightarrow F(P_A))}=\frac{F(P_0)}{ker(F(P_0)\rightarrow F(A))}=Im(F(P_0)\rightarrow F(A))=F(A)
			\end{gather*}
			\subsection{Universal $\delta$-Functority for $L_*F(-)$}
			The aim is to show $L_*F$ form a universal homological $\delta$-functor. 
			
			\begin{lemma}
				The objects $L_i F(A)$ of $\mathcal{B}$ are well defined up to isomorphism. That is: If $Q\rightarrow A$ is a second projective resolution, there exists a canonical isomorphism
				\begin{gather*}
					L_i F(A)=H_i(F(P))\stackrel{\cong}{\rightarrow} H_i(F(Q))
				\end{gather*}
				In particular, a different choice of projective resolution would yield new functor $\hat{L}_i F$, which are naturally isomorphic to the functors $L_iF$.
				\end{lemma}
				This implies, that the value of $L_iF(A)$
				\begin{proof}
					Apply comparision theorem:
					\begin{center}
						\begin{tikzcd}
							内容...
						\end{tikzcd}
					\end{center}
				\end{proof}
			\begin{corollary}
				If $A$ is projective, then $L_i F(A)=0,i\not=0$.
			\end{corollary}
			\begin{proof}
				Consider the projective resolution of $0\rightarrow P\stackrel{id}{\rightarrow} P\rightarrow 0$, then$L_iF(A)=H_i(F(P))=0$
			\end{proof}
			\subsection{$F$-Acyclic Functor}
			\begin{Def}
				An object $Q$ is called \textbf{$F$-Acyclic} if 
				\begin{gather*}
					L_iF(Q)=0,i\not=0
				\end{gather*}
			\end{Def}
			\subsection{$L_iF$ is Universal $\delta$-Functor}
			\begin{lemma}
				If $f:A'\rightarrow A$ is any map in $\mathcal{A}$, there exists a natural map $L_iF(f):L_iF(A')\rightarrow L_iF(A)$. 
			\end{lemma}
			\begin{proof}
				内容...
			\end{proof}
			\begin{proposition}
				内容...
			\end{proposition}
			\begin{proof}
				内容...
			\end{proof}
			\begin{theorem}
				Each $L_iF$ is an additive functor from $\mathcal{A}$ to $\mathcal{B}$
			\end{theorem}
			\begin{proof}
				内容...
			\end{proof}
			
			\begin{proposition}
				(Preversing Derived Functor)
			\end{proposition}
			\begin{proof}
				内容...
			\end{proof}
			
			\begin{theorem}
				The derived funcotr $L_*F$ form a homological $\delta$-functor.
			\end{theorem}
			\begin{proof}
				内容...
			\end{proof}
			\subsubsection{Dimension Shifting}
			\begin{proposition}
				内容...
			\end{proposition}
			\begin{proof}
				内容...
			\end{proof}
			\subsubsection{Homology,Cohomology}
			\begin{proposition}
				The functors $H_*:Ch_{\geq 0}(\mathcal{A})\rightarrow \mathcal{A}$ and $H^*:Ch^{\geq 0}(\mathcal{A})\rightarrow \mathcal{A}$ are universal $\delta$-functors.
			\end{proposition}
			\begin{proof}
				内容...
			\end{proof}
			\subsubsection{Effaceable Functors}
			\begin{Def}
				\begin{itemize}
					\item An additive funcotr $F:\mathcal{A}\rightarrow \mathcal{B}$ is \textbf{effaceable} if for each object $A$ of $\mathcal{A}$ there is a monomorphism $u:A\rightarrow I$ such that $F(u)=0$;
					\item Conversely, \textbf{coeffaceable} iff there exists surjection $u:P\rightarrow A$ such that $F(u)=0$.
				\end{itemize}
			\end{Def}
			\begin{proposition}
				\begin{itemize}
					\item If $T_*$ is a homological $\delta$-functor such that $T_n$ is coeffaceable except $T_0$, then $T_*$ is universal;
					\item If $T^*$ is a cohomological $\delta$-funcotr such that $T_n$ os effacable, then $T^*$ is universal. 
				\end{itemize}
			\end{proposition}
			\begin{proof}
				内容...
			\end{proof}
		\section{Right Derived Functors}
			\subsection{Right Derived Functors}
				\begin{Def}
					\begin{gather*}
						R^iF(A):=H^i(F(I))
					\end{gather*}
				\end{Def}
				\begin{proposition}
					(Right-Derived Functor is Dual of Left-Derived Functor)
				\end{proposition}
				\begin{proof}
					内容...
				\end{proof}
			\subsection{Ext Functor}
			\subsubsection{Ext Functor on $R$-modules}
				\begin{Def}
					Define:
					\begin{gather*}
						Ext^i_R(A,B):=R^iHom_R(A,-)(B)
					\end{gather*}
				\end{Def}
				
				\begin{proposition}
					(Criteria for Injective Modules) The followings are equivalent:
					\begin{itemize}
						\item $B$ is injective $R$-module;
						\item $Hom(-,B)$ is exact;
						\item $Ext_R^i(A,B)$ vanishes for all $i\not=0$ and all $A$; $B$ is $Hom_R(-,B)$ acyclic for all $A$.
						\item $Ext^1_R(A,B)$ vanishes for all $A$.
					\end{itemize}
				\end{proposition}
				\begin{proof}
					\begin{itemize}
						\item 
						\item 
						\item 
						\item 
					\end{itemize}
				\end{proof}
				
				\begin{proposition}
					(Criteria for Projective Modules)The followings are equivalent:
					\begin{itemize}
						\item $A$ is projective;
						\item $Hom_R(A,-)$ is exact;
						\item $Ext_R^i(A,B)$ vanishes for all $i\not=0$ and all $B$;
						\item $Ext^1_R(A,B)$ vanishes for all $B$.
					\end{itemize}
				\end{proposition}
				\begin{proof}
					内容...
				\end{proof}
			\subsubsection{Ext Functor in Abelian Categories}
				\begin{Def}
					\begin{itemize}
						\item 
						\item 
					\end{itemize}
				\end{Def}
				\begin{lemma}
					The two concepts above are equivalent in abelian category.
				\end{lemma}
				\begin{proof}
					内容...
				\end{proof}
				\begin{example}
					内容...
				\end{example}
			\subsubsection{Ext Functor in Derived Category}
				\begin{Def}
					Consider the \underline{derived category $D(\mathcal{A})$}. Define
					\begin{gather*}
						Ext^p(X,A):=Hom_{D(A)}(X,A[p])
					\end{gather*}
				\end{Def}
		\section{Adjoint Functors and Left/Right Exactness}
			\begin{theorem}
				Suppose $L:\mathcal{A}\rightarrow \mathcal{B}$ and $R:\mathcal{B}\rightarrow \mathcal{A}$ be adjoint pair. Then, there exists a natural isomorphism
				\begin{gather*}
					Hom(L(A),B)\stackrel{\cong}{\rightarrow} Hom(A,R(B))
				\end{gather*}
			\end{theorem}
			\begin{proof}
				内容...
			\end{proof}
			
			\begin{corollary}
				\begin{itemize}
				\item $\otimes_R B:mod-R\rightarrow Ab$ is left adjoint to $Hom_{Ab}(B,-)$, so $-\otimes_R B$ is right exact;
				\item More generally, if $B$ is an $R-S$ bimodule, then $\otimes_RB$ takes $mod-R$ to $mod-S$ and is a left adjoint, so it is right exact.
				\end{itemize}
			\end{corollary}
			\begin{proof}
				内容...
			\end{proof}
			
			\begin{proposition}
				内容...
			\end{proposition}
			\begin{proof}
				内容...
			\end{proof}
			\subsection{Tor Functor}
			\begin{Def}
				Define: For functor $T(A)=A\otimes_R B$, define the abelian group
				\begin{gather*}
					Tor^R_n(A,B):=(L_n T)(A)
				\end{gather*}
				
				In particular, $Tor_0^R(A,B)=A\otimes_R B$. 
				
				And when $A$
			\end{Def}
			\subsection{Adjoint Limit and Colimits}
		\section{Balancing Tor and Ext}
			Recall: Tensor Complex $(A\otimes B,d\otimes 1,(-1)^p\otimes d)$, take the notation by
			\begin{gather*}
				P\otimes_R Q\qquad Tot^\oplus(P\otimes_R Q)
			\end{gather*}
			
			\subsection{Balancing Tor}
			\begin{theorem} 
				For $R$-module $A,B$ we have:
				\begin{gather*}
					L_n(A\otimes_R -)(B)=L_n(\otimes_R B)(A)=Tor^R_n(A,B);n\in \NN
				\end{gather*}
			\end{theorem}
			\begin{proof}
				内容...
			\end{proof}
			
			\begin{corollary}
				We have:
				\begin{gather*}
					Tor^R_n(\oplus_{i\in I} A_i,B)\cong \oplus_{i\in I}Tor_n^R(A_i,B)\\
					Tor^R_n(A,\oplus_{i\in I} B_i)\cong \oplus_{i\in I} Tor_n^R(A,B_i)
				\end{gather*}
			\end{corollary}
			\begin{proof}
				内容...
			\end{proof}
			
			A question is: When can we state:
			\begin{gather*}
				Tor^R_n(\prod A_i,B)\cong \prod Tor^R_n(A_i,B)
			\end{gather*}
			
			\begin{itemize}
				\item $n=0$: $Tor^R_0(\prod A_i,B) \Leftrightarrow$
				\item $n\geq 1$: $\Leftrightarrow$ $R$ is right-coherent ring, and $B$ is finitely generated module.
			\end{itemize}
			
			\subsection{Balancing Ext}
			\begin{theorem}
				For $R$-module $A,B$ we have
				\begin{gather*}
					Ext^n_R(A,B)=R^n Hom_R(A,-)(B)=R^nHom_R(-,B)(A)
				\end{gather*}
			\end{theorem}
			\begin{proof}
				内容...
			\end{proof}
			
			\begin{corollary}
				We have
				\begin{gather*}
					Ext^n_R(A,\prod_{j\in J} B_j)=\prod_{j\in J} Ext_R^n(A,B_j)\\
					Ext^n_R(\oplus_{i\in I} A_i,B)=\prod_{i\in I} Ext^n_R(A_i,B)
				\end{gather*}
			\end{corollary}
			\begin{proof}
				内容...
			\end{proof}
		More discussions in next chapter.
	\chapter{Tor and Ext}
		\section{Tor for Abelian Groups}
			\begin{proposition}
				Suppose $A,B$ are abelian groups. Then we have
				\begin{gather*}
					Tor_n^\mathbb{Z}(A,B)=\begin{cases}
						\mathbb{Z}&n=1\\
						0&n>1
					\end{cases}
				\end{gather*}
			\end{proposition}
			\begin{proof}
				内容...
			\end{proof}
			
			\begin{theorem}
				(Acyclic Assembly Lemma) Let $C$ be a double complex in $mod-R$, then 
				\begin{itemize}
					\item 
					\item 
				\end{itemize}
			\end{theorem}
			\begin{proof}
				内容...
			\end{proof}
			
			\begin{Def}
				($Hom$-Cochain Complex)
			\end{Def}
			
			\begin{theorem}
				For every pair of $R$-modules $A$ and $B$, and all $n$,
				\begin{gather*}
					Ext_R^n(A,B)=R^nHom_R(A,-)(B)=R^nHom_R(-,B)(A)
				\end{gather*}
			\end{theorem}
			\begin{proof}
				内容...
			\end{proof}
			
			\begin{Def}
				内容...
			\end{Def}
		\section{Tor and Flatness}
			\begin{Def}
				Recall:
				\begin{itemize}
					\item 
					Consider $B\in R-Mod$. Then, $B$ is \textbf{flat} iff $-\otimes B$ is exact.
					
					\item Dual to define $A\in Mod-R$. Then, $A$ is \textbf{flat} iff $A\otimes -$ is exact.
				\end{itemize}
			\end{Def}
			\subsubsection{Pontrjagin Dual}
			\begin{Def}
				The Pontrjagin dual $B^*$ of a left-module $B$ is the right $R$-module 
				\begin{gather*}
					B\in Hom_{Ab}(B,\QQ/\ZZ)
				\end{gather*}
				and element $r\in R$ acts via $(fr)(b)=f(rb)$. So $B^*$ is a right $R$-module.
			\end{Def}
			
			\begin{theorem}
				For any $R-moudle$ we have
				\begin{gather*}
					Ext^n_R(A,B^\vee)\cong (Tor_n^R(A,B))^\vee
				\end{gather*}
			\end{theorem}
			\begin{proof}
				内容...
			\end{proof}
		\section{Ext for Nice Rings}
			\begin{lemma}
				Let $A,B$ be abelian groups, then we have
				\begin{gather*}
					Ext^n_{Ab}(A,B)=\begin{cases}
						\\
						0&n>1
					\end{cases}
				\end{gather*}
			\end{lemma}
			\begin{proof}
				内容...
			\end{proof}
		\section{Ext and Extensions}
		\section{Derived Functor and Inverse Limits}
			\textbf{Notion} In this section, we take
			\begin{itemize}
				\item $\mathcal{A}={\mathcal Ab}$;
				\item $I$ is the poset $\dots\rightarrow 2\rightarrow 1\rightarrow 0$;
			\end{itemize}
			\begin{Def}
				(Tower of Abelian Groups) Define
				\begin{gather*}
					\{A_i\}:\dots\rightarrow A_2\rightarrow A_1\rightarrow A_0
				\end{gather*}
			\end{Def}
			
			\subsection{Mittag-Leffers System}
			\begin{Def}
				(Mittag-Leffers System) Let $I$ be a directed set, $(A_i,\varphi_{ij})$ be an inverse system in $\mathcal Sets$. 
			\end{Def}
			\subsubsection{Limit}
			\begin{lemma}
				(Existence)
			\end{lemma}
			\begin{proof}
				内容...
			\end{proof}
			\begin{lemma}
				(Exactness) Let 
				\begin{gather*}
					内容...
				\end{gather*}
			\end{lemma}
		\section{Universal Coefficient Theorem}
			\begin{theorem}
				(K\"{u}nnenth) 
			\end{theorem}
			\begin{proof}
				内容...
			\end{proof}
			
			\begin{theorem}
				(UCT for Homology)
			\end{theorem}
			\begin{proof}
				内容...
			\end{proof}
			
			\begin{theorem}
				(UCT for Cohomology)
			\end{theorem}
			\begin{proof}
				内容...
			\end{proof}
			
			\subsection{Topology}
			\begin{theorem}
				(UCT in Topology)
			\end{theorem}
			\begin{proof}
				内容...
			\end{proof}
	\chapter{Homological Dimension}
		\section{Dimension}
			Define:
			\begin{Def}
				\begin{itemize}
					\item \textbf{The Projective Dimension}: 
					\item \textbf{The Injective Dimension}:
					\item \textbf{The Flat Dimension}
				\end{itemize}
			\end{Def}
			
			\begin{theorem}
				(Global Dimension Theorem) The following numbers are the same for any ring $R$:
				\begin{itemize}
					\item 
					\item 
					\item 
					\item 
				\end{itemize}
			\end{theorem}
			\begin{proof}
				内容...
			\end{proof}
			\begin{Def}
				内容...
			\end{Def}
			
			\subsubsection{Noetherian Ring}
			\subsection{Rings of Small Dimension}
		\section{Local Rings}
		\section{Konzul Complexes}
			\begin{Def}
				(Konzul Complex) Take the following notion
				\begin{gather*}
					内容...
				\end{gather*}
			\end{Def}
		\section{Local Cohomology}
		\section{Appendix: Hilbert-Serre Theorem}
			Our main reference:
			\begin{itemize}
				\item https://faculty.ustc.edu.cn/yqliang/files/myarticle/Serrethmonnoetherianregularlocalring.pdf
			\end{itemize}
			
			\begin{theorem}
				(Serre)
			\end{theorem}
			\begin{proof}
				内容...
			\end{proof}
	\chapter{Spectral Sequence}
		\section{Introductions}
			\begin{itemize}
				\item (Leray) Sheaf, Spectral Sequence;
				\item (Koszul,1945) Algebraic;
			\end{itemize}
			\subsection{Spectral Sequence}
			\begin{Def}
				Let $\mathcal{A}$ be an abelian category
				\begin{itemize}
					\item A \textbf{spectral sequence} in $\mathcal{A}$ is given by a system
					\begin{gather*}
						(E_r,d_r)_{r\geq 1}
					\end{gather*}
					where: Each $E_r$ is an object of $\mathcal{A}$, each $d_r:E_r\rightarrow E_r$ is a morphism such that
					\begin{gather*}
						d_r\circ d_r=0\qquad E_{r+1}=ker(d_r)/Im(d_r)
					\end{gather*}
					\item A morphism of spectral sequences
					\begin{gather*}
						f:(E_r,d_r)_{r\geq 1}\rightarrow (E_r',d_r')_{r\geq 1}
					\end{gather*}
					is given by a family of morphisms $(f_r)_{r\geq 1}$ such that
					\begin{gather*}
						f_r\circ d_r=d_r'\circ f_r
					\end{gather*}
					And $f_{r+1}:E_{r+1}\rightarrow E_{r+1}'$ is indentified by the morphism induced by $f_r:E_r\rightarrow E_r'$
					\begin{center}
						\begin{tikzcd}
							Im(d_r)\arrow[r,"\subseteq "]\arrow[d,"f_r|Im(d_r)"]&ker(d_r)\arrow[d,"f|ker(d_r)"]\arrow[r]&E_{r+1}\arrow[d,dashrightarrow]\\
							Im(d_r')\arrow[r,"\subseteq"]&ker(d_r')\arrow[r]&E_{r+1}'
						\end{tikzcd}
					\end{center}
				\end{itemize}
			\end{Def}
			\textbf{Remark} Sometimes denoted as $(E_r,d_r)_{r\geq 0}$.We loose the defition: There exists isomorphism $E_{r}\rightarrow ker(d_r)/Im(d_r)$.
			\begin{Def}
				内容...
			\end{Def}
			\subsubsection{Spectral Sequence}
				\begin{Def}
					Given a double complex $(E_r,d_r)_{r\geq 1}$, define
					\begin{gather*}
						0=B_1\subseteq B_2\subseteq \dots\subseteq Z_r\subseteq \dots\subseteq Z_2\subseteq Z_1=E_1
					\end{gather*}
					By following procedure:
					\begin{gather*}
						B_{r+1}=Im(d_r)\qquad Z_{r+1}=Ker(d_r)
					\end{gather*}
					Then, it is clear that $d_2:Z_2/B_2\rightarrow Z_2/B_2$. 
				\end{Def}
			\subsubsection{Limit of Spectral Sequence}
			\begin{Def}
				Let $\mathcal{A}$ be an abelian category, $(E_r,d_r)_{r\geq 1}$ be a spectral sequence.
				\begin{itemize}
					\item If the subobjects $Z_\infty=\cap Z_r$ and $B_\infty=\cup B_r$ of $E_1$ exist. Then, we define the \textbf{limit of spectral sequence} to be the object
					\begin{gather*}
						Z_\infty=\cap Z_r\qquad B_\infty=\cup B_r
					\end{gather*}
					\item We say that the spectral sequence degenerates at $E_r$ if the differentials $d_r,d_{r+1},\dots$ are all zero.
				\end{itemize}
			\end{Def}
			If the spectral sequence degenerates at $E_r$,then we have
			\begin{gather*}
				E_r=E_{r+1}=\dots=E_\infty
			\end{gather*}
			Also, almost any abelian category we
			will encounter has countable sums and intersections.
			
			\subsection{Exact Couple}
			\begin{Def}
				Let $\mathcal{A}$ be an exact category:
				\begin{itemize}
					\item An \textbf{exact couple} is a datum $(A,E,\alpha,f,g)$ where $A,E$ are objects of $\mathcal{A}$ and $\alpha,f,g$ are morphisms as in following diagram
					\begin{center}
						\begin{tikzcd}
							A\arrow[rr,"\alpha"]&&A\arrow[ld,"g"]\\
							&E
						\end{tikzcd}
					\end{center}
					\item 
				\end{itemize}
			\end{Def}
			\begin{lemma}
				内容...
			\end{lemma}
			\begin{proof}
				内容...
			\end{proof}
			\subsection{Differential Object}
			\begin{Def}
				Let $\mathcal{A}$ be an abeian category.
				
				A \textbf{differential object} of $\mathcal{A}$ is a pair $(A,d)$ consisting of
			\end{Def}
			\subsection{Filtered Differential Object}
			\subsection{Filtered Complexes}
			\subsection{Double Complexes}
			\subsection{Abelian Groups}
		
	\chapter{Simplicial Methods in Homological Algebra}
		\section{Simplicial Objects}
		\section{Operations on Simplicial Objects}
	\chapter{Hochshild and Cyclic Homology}
	\chapter{Derived Category}
		\section{Triangulated Category}
			\begin{Def}
				Define: Let $\mathcal{D}$ be an additive category, let
				\begin{gather*}
					[1]:\mathcal{D}\rightarrow \mathcal{D}\qquad E\mapsto E[1]
				\end{gather*}
				be an additive functor which is an auto-equivalence of $\mathcal{D}$:
				\begin{itemize}
					\item A \textbf{triangle} is sextuple
					\begin{gather*}
						(X,Y,Z,f,g,h)
					\end{gather*}
					where
					\begin{itemize}
						\item $X,Y,Z\in Ob(\mathcal{D})$;
						\item $f:X\rightarrow Y;g:Y\rightarrow Z;h:Z\rightarrow X[1]$ are morphisms of $\mathcal{D}$;
					\end{itemize}
					\item A morphism of triangles
					\begin{gather*}
						(X,Y,Z,f,g,h)\rightarrow (X',Y',Z',f',g',h')
					\end{gather*}
					with morphisms
					\begin{gather*}
						\begin{cases}
							a:X\rightarrow X'\\
							b:Y\rightarrow Y'\\
							c:Z\rightarrow Z'
						\end{cases}
					\end{gather*}
					of $\mathcal{D}$ such that
					\begin{gather*}
						\begin{cases}
							b\circ f=f'\circ a\\
							c\circ g=g'\circ b\\
							a[1]\circ h=h'\circ c
						\end{cases}
					\end{gather*}
				\end{itemize}
				With visualization
				\begin{center}
					\begin{tikzcd}
						X&Y&Z&X[1]\\
						X'&Y'&Z'&X'[1]
					\end{tikzcd}
				\end{center}
			\end{Def}
			
			\subsection{Triangulated Category}
			\begin{Def}
				Axioms:
				\begin{itemize}
					\item TR1: Any triangle category isomorphic to a distinguished triangle is a distinguished triangle. 
					
					Any triangle of the form $(X,X,0,id,0,0)$ is distinguished.
					
					For any morphism $f:X\rightarrow Y$ in $\mathcal{D}$, there exists a \textbf{distinguished triangle} of the form
					\begin{gather*}
						(X,Y,Z,f,g,h)
					\end{gather*}
					\item TR2: The triangle $(X,Y,Z,f,g,h)$ is distinguished iff $(Y,Z,X[1],g,h,-f[1])$ is
					\item TR3: Given the diagram
					\begin{center}
						\begin{tikzcd}
							X\arrow[r,"f"]&Y\arrow[r,"g"]&Z\arrow[r,"h"]&X[1]\\
							X'\arrow[r]&Y'\arrow[r]&Z'\arrow[r]&X‘[1]
						\end{tikzcd}
					\end{center}
					\item TR4: Given objects $X,Y,Z\in Ob(\mathcal{D})$, and morphisms $f:X\rightarrow Y,g:Y\rightarrow Z$ and distinguished triple
					\begin{gather*}
						(X,Y,Q_1,f,p_1,d_1)\quad (X,Z,g\circ f,p_2,d_2)\quad (Y,Z,Q_2,g,p_3,d_3)
					\end{gather*}
					There exists morphisms 
					\begin{gather*}
						a:Q_1\rightarrow Q_2\qquad b:Q_2\rightarrow Q_3
					\end{gather*}
					such that:
					\begin{itemize}
						\item $(Q_1,Q_2,Q_3,)$
						\item 
						\item 
					\end{itemize}
				\end{itemize}
				We will call:
				\begin{itemize}
					\item 
					\item 
				\end{itemize}
			\end{Def}
			\textbf{Remark}
			\subsubsection{Homological Functor}
			\begin{Def}
				\begin{itemize}
					\item Let $\mathcal{D}$ be a pre-triangulated category, $\mathcal{A}$ be an abelian category.
					
					An additive category $H:\mathcal{D}\rightarrow \mathcal{A}$ is \textbf{homological} iff for every distinguished  triangle $(X,Y,Z,f,g,h)$ the sequence
					\begin{gather*}
						H(X)\rightarrow H(Y)\rightarrow H(Z)
					\end{gather*}
					is exact in the abelian category $\mathcal{A}$
					\item Conversely, $H:\mathcal{D}^{op}\rightarrow \mathcal{A}$ if the corresponding functor
					\begin{gather*}
						H^{op}:\mathcal{D}\rightarrow \mathcal{A}^{op}
					\end{gather*}
					is homological.
				\end{itemize}.
			\end{Def}
			\textbf{Remark} If $H:\mathcal{D}\rightarrow \mathcal{A}$, we write
			\begin{gather*}
				H^n(X):=H(X(n))
			\end{gather*}
				\begin{Def}
					By TR2, a distinguished triangle $(X,Y,Z,f,g,h)$ determines a long exact sequence
					\begin{gather*}
						H^{-1}(Z)\stackrel{above}{\rightarrow} H^0(X)\stackrel{above}{\rightarrow} H^0(Y)\stackrel{above}{\rightarrow} H^0(Z)\stackrel{H(h)}{\rightarrow} H^1(X)
					\end{gather*}
				\end{Def}
				\subsubsection{$\delta$-Functor}
				\begin{Def}
					Let $\mathcal{A}$ be an abelian category. Let $\mathcal{D}$ be a triangulated category. 
					
					A $\delta$-functor from $\mathcal{A}$ to $\mathcal{D}$ is given by a functor
					\begin{gather*}
						G:\mathcal{A}\rightarrow \mathcal{D}
					\end{gather*}
				\end{Def}
			\subsection{Elementray Results}
				\begin{lemma}
					Let $\mathcal{D}$ be a pre-triangulated category, let $(X,Y,Z,f,g,h)$ be distinguished triangle, then
					\begin{gather*}
						g\circ f=0\quad h\circ g=0\quad f[1]\circ h=0
					\end{gather*}
				\end{lemma}
				\begin{proof}
					内容...
				\end{proof}
				
				\begin{lemma}
					Let $\mathcal{D}$ be a pre-triangulated category. For any object $W$ of $\mathcal{D}$, the functor $Hom_\mathcal{D}(W,-)$ is homological, and the functor $Hom_{\mathcal{D}}(-,W)$ is cohomological.
				\end{lemma}
				\begin{proof}
					内容...
				\end{proof}
				
				\begin{lemma}
					Let $\mathcal{D}$ be a pretriangulated category, and
					\begin{gather*}
						(a,b,c):(X,Y,Z,f,g,h)\rightarrow 
					\end{gather*}
				\end{lemma}
				\begin{proof}
					内容...
				\end{proof}
				
				\begin{proposition}
					Let $\mathcal{D}$ be \underline{trianglatedd category}. Any commutative category
					\begin{center}
						\begin{tikzcd}
							X&Y\\
							X'&Y'
						\end{tikzcd}
					\end{center}
					can be extended to the diagram:
					\begin{center}
						\begin{tikzcd}
							X&Y&Z&X[1]\\
							X'&Y'&Z'&X'[1]\\
							X''&Y''&Z‘’&X''[1]\\
							X[1]&Y[1]&Z[1]&X[2]
						\end{tikzcd}
					\end{center}
					where: All square commutes, except for the lower right squrae which is anticommutative.
					
					Moreover, each of column and rows are distinguished triangles.
					
					Finally, the morphisms on the bottom row/right column are obtained from the morphisms of the top row/left column applying $[1]$.
				\end{proposition}
				\begin{proof}
					内容...
				\end{proof}
			\subsection{Constructions}
			\subsubsection{Triangulated Functor}
			\subsubsection{Localization of Triangulated Functor}
				\begin{Def}
					Let: $\mathcal{D}$ be a pre-triangulated category, we say a multiplicative system $S$ is \textbf{compatible with triangulated stucture} if the following two conditions hold:
					\begin{itemize}
						\item MS5: For a morphism $f$ of $\mathcal{D}$ we have
						\begin{gather*}
							f\in S\Leftrightarrow f[1]\in S
						\end{gather*}
						\item Given a solid commutative square
						\begin{center}
							\begin{tikzcd}
								X\arrow[r]\arrow[d,"s"]&Y\arrow[r]\arrow[d,"{s'}"]&Z\arrow[d,dashrightarrow,"{\exists s''}"]\arrow[r]&X[1]\arrow[d,"{s[1]}"]\\
								X'\arrow[r]&Y'\arrow[r]&Z'\arrow[r]&X'[1]
							\end{tikzcd}
						\end{center}
						whose rows are distinguished triangles with $s,s'\in S$, there exists a morphism
						\begin{gather*}
							s'':Z\rightarrow Z'
						\end{gather*}
						in $S$ such that $(s,s',s'')$ is a morphism of triangles.
					\end{itemize}
				\end{Def}
				\begin{lemma}
					内容...
				\end{lemma}
			\subsubsection{Quotient of Triangulated Category}
			\begin{lemma}
				
			\end{lemma}
		\section{The Homotopy Category}
			\begin{Def}
				Let $\mathcal{A}$ be additive category:
				\begin{itemize}
					\item We set $Comp(\mathcal{A})=CoCh(\mathcal{A})$ be the \textbf{category of (cochain) complexes}.
					\item A complex $K^\bullet$ is \textbf{bounded below} if:
					\begin{gather*}
						内容...
					\end{gather*}
					\item 
					\item 
					\item 
					\item 
					\item 
				\end{itemize}
			\end{Def}
		\section{Derived Category}
			\subsection{Filtered Derived Category}
		\chapter{Appendix}
			\section{Eckmann-Hilton}
				\begin{theorem}
					Suppose $X$ be a structure, with two binary operations $\circ,\star$ which satisfies
					\begin{itemize}
						\item (Identity) $x\circ e=e\circ x=x;x\star e=e\star x=x$;
						\item (Eckmann-Hilton) For $a,b,c,d\in X$,
						\begin{gather*}
							(a\circ b)\star(c\circ d)=(a\star c)\circ (b\star d)
						\end{gather*}
						
						Then:
						\begin{itemize}
							\item $a\circ b=a\star b$;
							\item The operation is commutative;
							\item The operation is associative.
						\end{itemize}
					\end{itemize}
				\end{theorem}
				\begin{proof}
					Notice
					\begin{itemize}
						\item 
						\item 
					\end{itemize}
				\end{proof}
	\end{document}
	